\documentclass[11pt,a4paper]{article}
\usepackage[utf8]{inputenc}\usepackage[T1]{fontenc}\usepackage[italian,english]{babel}
\usepackage{amsmath,amssymb,amsthm}\usepackage[margin=2.2cm]{geometry}\usepackage{longtable,booktabs,array,calc}
\usepackage{listings,xcolor}\usepackage{hyperref}\usepackage{url}
\providecommand{\tightlist}{\setlength{\itemsep}{0pt}\setlength{\parskip}{0pt}}
\lstset{basicstyle=\ttfamily\scriptsize,breaklines=true,frame=single,captionpos=b,columns=fullflexible,keepspaces=true,inputencoding=utf8,extendedchars=true}
\title{Proof Pursuit --- rapporto completo\\ \large prove, decisioni degli agenti, codice verificato, letteratura}
\author{Team Proof Pursuit (Researcher: T. Tumini; Referee: gabundos; Creative: M. Renzi; gio3r3gio)}
\date{\today}
\begin{document}\maketitle\tableofcontents
\section*{Come leggere questo rapporto}
Per ogni cella: la richiesta ufficiale, lo stato dichiarato dal team (mai ``risolto'' senza revisione umana), la consegna
con la prova, poi la traccia delle decisioni degli agenti (Researcher: famiglia e motivo; Referee: verdetto ed errore
fatale; Creative: quando e perché è intervenuto) e il codice su cui il tentativo poggia con l'esito della riesecuzione
fidata. DIMOSTRATO, CITATO, VERIFICATO SU INTERVALLO FINITO e CONGETTURA restano distinti. L'architettura del sistema
è descritta nel README del repository.
\section{Problema 1 — Angles between lines (Fejes Toth)}
\subsection{Enunciato ufficiale: preambolo e regole di consegna}
\subsubsection{Problema 1 --- Angles between
lines}\label{problema-1-angles-between-lines}

\emph{(Enunciato formattato in LaTeX, contenuto fedele all'originale;
ricevuto a parti il 2026-09-26.)}

\paragraph{Preambolo}\label{preambolo}

\textbf{Angles between lines.} How large can the sum of pairwise angles
between \(N\) lines through the origin be? Fejes Tóth conjectured in
1959 that orthogonal lines win.

A line here always means a line through the origin of \(\mathbb{R}^d\).
The angle between two lines \(\ell, \ell'\) is the acute (non-obtuse)
angle \(\theta(\ell,\ell') \in [0,\pi/2]\) between them. If
\(\ell, \ell'\) are spanned by unit vectors \(x, x'\), then \[
\theta(\ell,\ell') = \arccos\,\lvert \langle x, x' \rangle \rvert .
\] For lines \(\ell_1, \ldots, \ell_N\) in \(\mathbb{R}^d\) (repetitions
allowed), write \[
S(\ell_1,\ldots,\ell_N) = \sum_{1 \le i < j \le N} \theta(\ell_i,\ell_j).
\] The question is how large \(S\) can be. In 1959 L. Fejes Tóth
conjectured that \(S\) is maximised by taking \(d\) mutually orthogonal
lines, each used either \(\lfloor N/d \rfloor\) or \(\lceil N/d \rceil\)
times. When \(N = d + k\) with \(0 \le k \le d\), that configuration
(the \(d\) coordinate axes, \(k\) of them used twice) has \[
S = \left( \binom{N}{2} - k \right) \cdot \frac{\pi}{2},
\] because exactly \(k\) pairs of lines coincide and every other pair is
orthogonal.

Each cell below asks you to prove this bound, or a special case of it.
Unless a cell says otherwise, you need a complete proof. Citing a
published result for the statement you are asked to prove does not
count.

\textbf{What to hand in.} For each cell you attempt, hand in a written
proof. You may use a computer to explore. A proof may rely on a
computation only if you include the code, it runs in under 10 minutes on
a laptop, and it is rigorous: exact or interval arithmetic, or an
argument that bounds the numerical error. Say clearly which cells you
consider solved and which are partial.

\subsection{Cella 1 (1 punti) — stato: in corso}
\subparagraph{Parte 1 --- Lines in the
plane}\label{parte-1-lines-in-the-plane}

\textbf{Punteggio:} 1 point · \textbf{Valutazione:} Judged

We begin in the plane (\(d = 2\)), where the conjectured optimum splits
the \(N\) lines as evenly as possible between two perpendicular
directions. Let \(\ell_1, \ldots, \ell_N\) be lines in \(\mathbb{R}^2\).
Prove that \[
S(\ell_1, \ldots, \ell_N) \;\le\; \frac{\pi}{2} \cdot \left\lfloor \frac{N^2}{4} \right\rfloor .
\]

\textbf{Enunciato protetto dato agli agenti (state.json).} \texttt{Let l\_1..l\_N be lines through the origin in R\textasciicircum{}2. Prove S(l\_1..l\_N) = sum\_\{i\textless{}j\} theta(l\_i,l\_j) \textless{}= (pi/2) * floor(N\textasciicircum{}2/4).}

\textbf{Evidenza registrata in STATUS.md.} E1: max numerico = target per N≤10; prova via argomento di Crofton abbozzata (note.md, E), da scrivere e revisionare; definizione di \$S\$ confermata dal preambolo

\subsubsection{Consegna (prova)}
Nessuna bozza di consegna ancora scritta.

\subsubsection{Decisioni degli agenti}
\paragraph{Run \texttt{p1\_c1}}

\subsubsection{Codice su cui poggia il tentativo}
Nessun codice nei tentativi.

\subsection{Cella 2 (2 punti) — stato: in corso}
\subparagraph{Parte 2 --- An orthogonality
lemma}\label{parte-2-an-orthogonality-lemma}

\textbf{Punteggio:} 2 points · \textbf{Valutazione:} Judged

A statement about a chain of vectors, in which each vector may fail to
be orthogonal only to its immediate neighbours in the list. Let
\(m \ge 2\) and let \(x_1, \ldots, x_m\) be unit vectors in
\(\mathbb{R}^{m-1}\) with \(\langle x_i, x_j \rangle = 0\) whenever
\(|i - j| \ge 2\). Prove that \[
\sum_{i=1}^{m-1} \theta(x_i, x_{i+1}) \;\le\; (m-2) \cdot \frac{\pi}{2},
\] where \(\theta(x,y) = \arccos\,\lvert \langle x, y \rangle \rvert\).

\textbf{Enunciato protetto dato agli agenti (state.json).} \texttt{Let m \textgreater{}= 2 and x\_1..x\_m unit vectors in R\textasciicircum{}\{m-1\} with \textless{}x\_i,x\_j\textgreater{} = 0 whenever |i-j| \textgreater{}= 2. Prove sum\_\{i=1\}\textasciicircum{}\{m-1\} theta(x\_i,x\_\{i+1\}) \textless{}= (m-2) pi/2, where theta(x,y) = arccos|\textless{}x,y\textgreater{}|.}

\textbf{Evidenza registrata in STATUS.md.} prova per induzione abbozzata (note.md, F); casi m=2,3 verificati a mano

\subsubsection{Consegna (prova)}
Nessuna bozza di consegna ancora scritta.

\subsubsection{Decisioni degli agenti}
\paragraph{Run \texttt{p1\_c2}}

\subsubsection{Codice su cui poggia il tentativo}
Nessun codice nei tentativi.

\subsection{Cella 3 (3 punti) — stato: in analisi}
\subparagraph{\texorpdfstring{Parte 3 --- \(d+1\) lines in
\(\mathbb{R}^d\)}{Parte 3 --- d+1 lines in \textbackslash mathbb\{R\}\^{}d}}\label{parte-3-d1-lines-in-mathbbrd}

\textbf{Punteggio:} 3 points · \textbf{Valutazione:} Judged

The first case in every dimension: one line more than the dimension,
\(N = d+1\), where the conjectured optimum repeats exactly one of the
\(d\) coordinate axes. Let \(d \ge 1\). Prove that any \(d+1\) lines in
\(\mathbb{R}^d\) satisfy \[
S \;\le\; \left( \binom{d+1}{2} - 1 \right) \cdot \frac{\pi}{2}.
\]

\textbf{Enunciato protetto dato agli agenti (state.json).} \texttt{Let d \textgreater{}= 1. Prove that any d+1 lines through the origin in R\textasciicircum{}d satisfy S \textless{}= (C(d+1,2) - 1) pi/2.}

\textbf{Evidenza registrata in STATUS.md.} triage in note.md; dipende dal lemma P2

\subsubsection{Consegna (prova)}
Nessuna bozza di consegna ancora scritta.

\subsubsection{Decisioni degli agenti}
\paragraph{Run \texttt{p1\_c3}}

\subsubsection{Codice su cui poggia il tentativo}
Nessun codice nei tentativi.

\subsection{Cella 4 (5 punti) — stato: in analisi}
\subparagraph{\texorpdfstring{Parte 4 --- Five lines in
\(\mathbb{R}^3\), six in
\(\mathbb{R}^4\)}{Parte 4 --- Five lines in \textbackslash mathbb\{R\}\^{}3, six in \textbackslash mathbb\{R\}\^{}4}}\label{parte-4-five-lines-in-mathbbr3-six-in-mathbbr4}

\textbf{Punteggio:} 5 points · \textbf{Valutazione:} Judged

Two concrete instances of the next case, \(N = d+2\): five lines in
three dimensions and six lines in four. In each, the conjectured optimum
repeats two of the coordinate axes. Prove that any \(5\) lines in
\(\mathbb{R}^3\) satisfy \(S \le 4\pi\), and that any \(6\) lines in
\(\mathbb{R}^4\) satisfy \(S \le 13\pi/2\).

\textbf{Enunciato protetto dato agli agenti (state.json).} \texttt{Prove that any 5 lines through the origin in R\textasciicircum{}3 satisfy S \textless{}= 4 pi, and that any 6 lines in R\textasciicircum{}4 satisfy S \textless{}= 13 pi/2. A rigorous computation (exact/interval arithmetic, \textless{} 10 min, code included) is admissible.}

\textbf{Evidenza registrata in STATUS.md.} triage; possibile via P5 o calcolo rigoroso

\subsubsection{Consegna (prova)}
Nessuna bozza di consegna ancora scritta.

\subsubsection{Decisioni degli agenti}
\paragraph{Run \texttt{p1\_c4}}

\subsubsection{Codice su cui poggia il tentativo}
Nessun codice nei tentativi.

\subsection{Cella 5 (8 punti) — stato: in analisi}
\subparagraph{\texorpdfstring{Parte 5 --- \(d+2\) lines in
\(\mathbb{R}^d\)}{Parte 5 --- d+2 lines in \textbackslash mathbb\{R\}\^{}d}}\label{parte-5-d2-lines-in-mathbbrd}

\textbf{Punteggio:} 8 points · \textbf{Valutazione:} Judged

The case \(N = d+2\) in every dimension, with the same conjectured
optimum: the \(d\) coordinate axes, two of them repeated. Prove that for
every \(d \ge 2\), any \(d+2\) lines in \(\mathbb{R}^d\) satisfy \[
S \;\le\; \left( \binom{d+2}{2} - 2 \right) \frac{\pi}{2}.
\]

\textbf{Enunciato protetto dato agli agenti (state.json).} \texttt{Prove that for every d \textgreater{}= 2, any d+2 lines through the origin in R\textasciicircum{}d satisfy S \textless{}= (C(d+2,2) - 2) pi/2.}

\textbf{Evidenza registrata in STATUS.md.} triage

\subsubsection{Consegna (prova)}
Nessuna bozza di consegna ancora scritta.

\subsubsection{Decisioni degli agenti}
\paragraph{Run \texttt{p1\_c5}}

\subsubsection{Codice su cui poggia il tentativo}
Nessun codice nei tentativi.

\subsection{Cella 6 (13 punti) — stato: in analisi}
\subparagraph{\texorpdfstring{Parte 6 --- \(N\) lines in
\(\mathbb{R}^d\)}{Parte 6 --- N lines in \textbackslash mathbb\{R\}\^{}d}}\label{parte-6-n-lines-in-mathbbrd}

\textbf{Punteggio:} 13 points · \textbf{Valutazione:} Judged ·
\textbf{Open question}

Fejes Tóth's conjecture from the Setting, for every \(N\) and every
\(d\). For \(N\) lines in \(\mathbb{R}^d\), write \(N = qd + s\) with
\(0 \le s < d\), and let \[
M(N,d) = s \binom{q+1}{2} + (d-s) \binom{q}{2}.
\] Prove or disprove: every \(N\) lines in \(\mathbb{R}^d\) satisfy \[
S \;\le\; \left( \binom{N}{2} - M(N,d) \right) \frac{\pi}{2}.
\] This is the value attained by splitting the lines as evenly as
possible among \(d\) mutually orthogonal directions. Settling any
infinite family not already covered above counts as partial progress.

\emph{(Fine del problema 1. Punteggi: 1+2+3+5+8+13 = 32.)}

\textbf{Enunciato protetto dato agli agenti (state.json).} \texttt{OPEN. For N lines in R\textasciicircum{}d write N = qd+s, 0 \textless{}= s \textless{} d, M(N,d) = s C(q+1,2) + (d-s) C(q,2). Prove or disprove: S \textless{}= (C(N,2) - M(N,d)) pi/2. Settling any infinite family not covered by cells 1-5 counts as partial progress.}

\textbf{Evidenza registrata in STATUS.md.} triage; solo ricerca controesempi / famiglie infinite

\subsubsection{Consegna (prova)}
Nessuna bozza di consegna ancora scritta.

\subsubsection{Decisioni degli agenti}
\paragraph{Run \texttt{p1\_c6}}

\subsubsection{Codice su cui poggia il tentativo}
Nessun codice nei tentativi.

\subsection{Letteratura arXiv consultata (ricerca deterministica, abstract letti)}
\begin{itemize}
\item \cite{1801.07837v1} On the Fejes Tóth Problem about the Sum of Angles Between Lines — Dmitriy Bilyk, Ryan W Matzke (2018). \emph{Query}: \texttt{angles between lines}.
\item \cite{2007.08698v2} On Fejes Tóth's conjectured maximizer for the sum of angles between lines — Tongseok Lim, Robert J. McCann (2020). \emph{Query}: \texttt{angles between lines}.
\item \cite{1204.3850v1} Simple Agents Learn to Find Their Way: An Introduction on Mapping Polygons — Jérémie Chalopin, Shantanu Das, Yann Disser, Matúš Mihalák (2012). \emph{Query}: \texttt{angles between lines}.
\end{itemize}
\section{Problema 2 — Uphill paths on the hypercube}
\subsection{Enunciato ufficiale: preambolo e regole di consegna}
\subsubsection{Problema 2 --- Uphill paths on the
hypercube}\label{problema-2-uphill-paths-on-the-hypercube}

\emph{(Enunciato formattato in LaTeX, contenuto fedele all'originale;
ricevuto il 2026-09-26.)}

\paragraph{Preambolo}\label{preambolo}

\textbf{Uphill paths on the hypercube.} Label the vertices of the
hypercube to make as few uphill paths as possible. IMO 2022 settled the
grid; the cube is still open.

Let \(G\) be a finite simple graph on \(n\) vertices. A labelling of
\(G\) is a bijection \(f\) from \(V(G)\) to \(\{1,2,\ldots,n\}\). Fix a
labelling.

\begin{itemize}
\tightlist
\item
  A vertex \(v\) is a \emph{valley} if every neighbour \(w\) of \(v\)
  has \(f(w) > f(v)\). An isolated vertex counts as a valley.
\item
  An \emph{uphill path} is a sequence \((v_1, v_2, \ldots, v_k)\) with
  \(k \ge 1\) in which \(v_1\) is a valley, \(v_i\) and \(v_{i+1}\) are
  adjacent for each \(i\), and \(f(v_1) < f(v_2) < \cdots < f(v_k)\). A
  valley on its own (\(k=1\)) counts as an uphill path.
\end{itemize}

Let \(U(G)\) be the smallest possible number of uphill paths, taken over
all labellings of \(G\).

IMO 2022 Problem 6 asks for \(U\) of the \(n \times n\) grid graph. The
answer is \(2n^2 - 2n + 1\).

This column asks about the \(d\)-dimensional hypercube \(Q_d\). Its
vertex set is \(\{0,1\}^d\), and two vertices are adjacent when they
differ in exactly one coordinate. \(Q_d\) has \(2^d\) vertices and
\(d \cdot 2^{d-1}\) edges.

\textbf{What to hand in.} For C1 to C4, hand in the values, and for each
value an explicit labelling that attains it as a list of the \(2^d\)
vertices in increasing label order, written as 0/1 strings. These cells
are marked correct on the values alone, but a value with no labelling
behind it will not survive the later cells, which build on the
construction.

For C5, hand in either a labelling of \(Q_9\) in the same format, whose
uphill paths will be counted mechanically, or a proof of the lower
bound.

For C6, hand in all three of: the value; an explicit labelling that
attains it, in the same format; a proof that no labelling gives fewer
uphill paths.

A lower-bound proof may be computer-assisted. If it is, include the
code; it must run in under 10 minutes on a laptop, and you must explain
why the computation proves the bound. Say clearly which cells you
consider solved and which are partial.

\subsection{Cella 1 (1 punti) — stato: revisionato}
\subparagraph{\texorpdfstring{Parte 1 (C1) --- \(U(Q_3)\) and
\(U(Q_4)\)}{Parte 1 (C1) --- U(Q\_3) and U(Q\_4)}}\label{parte-1-c1-uq_3-and-uq_4}

\textbf{Punteggio:} 1 point · \textbf{Valutazione:} Checked instantly

The two smallest interesting cubes: \(Q_3\) has \(8\) vertices and
\(12\) edges, \(Q_4\) has \(16\) vertices and \(32\) edges. Determine
\(U(Q_3)\) and \(U(Q_4)\).

\textbf{Enunciato protetto dato agli agenti (state.json).} \texttt{Determine U(Q\_3) and U(Q\_4), each with an explicit optimal labelling (list of the 2\textasciicircum{}d vertices in increasing label order as 0/1 strings).}

\textbf{Evidenza registrata in STATUS.md.} \$U(Q\_3)=14\$, \$U(Q\_4)=34\$: prove a mano rilette passo per passo + verifiche esatte indipendenti; bozza in submission/parte-1.md

\subsubsection{Consegna (prova)}
\paragraph{Problema 2 --- Parte 1 (C1) --- bozza di
consegna}\label{problema-2-parte-1-c1-bozza-di-consegna}

\textbf{Stato dichiarato: RISOLTA (bozza revisionata).} Entrambe le
metà: \(U(Q_3)=14\) e \(U(Q_4)=34\), con valore, etichettatura esplicita
e prova completa a mano; calcoli esatti a conferma.

\subparagraph{1. Risultato e ambito}\label{risultato-e-ambito}

\(U(Q_3) = 14\) e \(U(Q_4) = 34\). Etichettatura ottima di \(Q_4\)
(ordine crescente di etichetta; due valli, \texttt{0000} e
\texttt{1111}):
\texttt{0000,\ 1111,\ 0111,\ 1101,\ 1110,\ 1011,\ 0001,\ 1000,\ 0010,\ 0100,\ 1001,\ 0110,\ 0011,\ 1010,\ 0101,\ 1100}.

Per \(Q_3\): Etichettatura ottima (ordine crescente di etichetta;
posizione \(i\) della stringa = coordinata \(i\); la convenzione è
irrilevante perché le permutazioni di coordinate sono automorfismi di
\(Q_3\)): \texttt{000,\ 100,\ 010,\ 110,\ 101,\ 011,\ 001,\ 111}.

\subparagraph{2. Dimostrazione}\label{dimostrazione}

\subparagraph{Result}\label{result}

\[U(Q_3) = 14.\]

An optimal labelling, as the list of the 8 vertices in increasing label
order (string position \(i\) = coordinate \(i\); the convention is
irrelevant for the count since \(\mathrm{Aut}(Q_3)\) contains all
coordinate permutations):

\[\texttt{000},\ \texttt{100},\ \texttt{010},\ \texttt{110},\ \texttt{101},\ \texttt{011},\ \texttt{001},\ \texttt{111}\]

i.e.~\(f(000)=1,\ f(100)=2,\ f(010)=3,\ f(110)=4,\ f(101)=5,\ f(011)=6,\ f(001)=7,\ f(111)=8\).

\subparagraph{Notation and the counting
identity}\label{notation-and-the-counting-identity}

For a labelling \(f\) of a graph \(G\), orient every edge from the
smaller to the larger label; this gives an acyclic orientation. Write
\(u \to v\) if \(uv\in E\) and \(f(u)<f(v)\), and
\(\deg^-(v)=\#\{u: u\to v\}\). Let \(p(v)\) be the number of uphill
paths whose last vertex is \(v\).

\textbf{Lemma 1 (recursion).}
\(p(v) = [\,v \text{ is a valley}\,] + \sum_{u\to v} p(u)\), and the
total number of uphill paths is \(\sum_v p(v)\).

\emph{Proof.} An uphill path ending at \(v\) has \(k=1\) (then it is
\((v)\) and \(v\) must be a valley; conversely a valley gives exactly
this one path) or \(k\ge2\); in the latter case deleting \(v\) gives an
uphill path ending at \(v_{k-1}\) with \(v_{k-1}\to v\), and conversely
appending \(v\) to any uphill path ending at some \(u\) with \(u\to v\)
gives an uphill path ending at \(v\) (the label condition \(f(u)<f(v)\)
is exactly \(u\to v\)). These correspondences are bijective, and the
sets for different \(u\) are disjoint (they differ in the penultimate
vertex). Every uphill path has a unique last vertex, hence the total is
\(\sum_v p(v)\). \(\square\)

\subparagraph{Upper bound: the labelling above has exactly 14 uphill
paths}\label{upper-bound-the-labelling-above-has-exactly-14-uphill-paths}

Neighbours in \(Q_3\): flip one coordinate. Compute \(p\) in increasing
label order using Lemma 1:

{\def\LTcaptype{none} % do not increment counter
\begin{longtable}[]{@{}llll@{}}
\toprule\noalign{}
vertex & label & lower neighbours & \(p\) \\
\midrule\noalign{}
\endhead
\bottomrule\noalign{}
\endlastfoot
000 & 1 & none (valley) & 1 \\
100 & 2 & 000 & 1 \\
010 & 3 & 000 & 1 \\
110 & 4 & 100, 010 & 1+1 = 2 \\
101 & 5 & 100 (001 has label 7, 111 has 8) & 1 \\
011 & 6 & 010 (001 has 7, 111 has 8) & 1 \\
001 & 7 & 000, 101, 011 & 1+1+1 = 3 \\
111 & 8 & 110, 101, 011 & 2+1+1 = 4 \\
\end{longtable}
}

Only 000 is a valley (every other vertex has a lower neighbour, as the
table shows). Total \(=1+1+1+2+1+1+3+4 = 14\). Hence \(U(Q_3)\le 14\).

\subparagraph{\texorpdfstring{Lower bound (hand proof): no labelling of
\(Q_3\) has \(\le 13\) uphill
paths}{Lower bound (hand proof): no labelling of Q\_3 has \textbackslash le 13 uphill paths}}\label{lower-bound-hand-proof-no-labelling-of-q_3-has-le-13-uphill-paths}

\textbf{Lemma 2.} For every labelling of any graph and every vertex
\(v\): \(p(v)\ge 1\), and \(p(v)\ge \deg^-(v)\).

\emph{Proof.} Starting from \(v\), repeatedly move to a neighbour with
smaller label while one exists; labels strictly decrease so this stops,
at a valley \(v_1\). Reversing the walk gives an uphill path ending at
\(v\), so \(p(v)\ge1\). Applying this to each \(u\) with \(u\to v\)
gives \(p(u)\ge1\), so by Lemma 1
\(p(v)\ge\sum_{u\to v}p(u)\ge\deg^-(v)\). \(\square\)

\textbf{Corollary 3.} Total
\(\ \ge \sum_v \max(1,\deg^-(v)) = \#\{\text{valleys}\} + \sum_{v \text{ not valley}}\deg^-(v) = \#\{\text{valleys}\} + |E|\),
since valleys are exactly the vertices with \(\deg^-=0\) and
\(\sum_v\deg^-(v)=|E|\) (each edge has exactly one head). For \(Q_3\),
\(|E|=12\) and the vertex with label 1 is always a valley, so the total
is \(\ge 13\).

\textbf{Proposition 4.} The total is never exactly 13 for \(Q_3\); hence
\(U(Q_3)\ge 14\).

\emph{Proof.} Suppose a labelling has total 13. Then every inequality in
Corollary 3 is an equality: (a) there is exactly one valley, and (b) for
every non-valley \(v\), \(p(v)=\deg^-(v)\),
i.e.~\(\sum_{u\to v}p(u)=\deg^-(v)\) with each \(p(u)\ge1\), so
\(p(u)=1\) for every \(u\to v\). By Lemma 2, \(p(u)=1\) forces
\(\deg^-(u)\le1\).

Let \(x\) be the vertex with label 8. All 3 neighbours of \(x\) are
lower, so \(\deg^-(x)=3\) and by (b) every \(a\in N(x)\) has
\(\deg^-(a)\le1\).

Let \(y\) be the vertex with label 7. At most one neighbour of \(y\)
(namely \(x\)) is higher, so \(\deg^-(y)\ge2\). If \(y\in N(x)\) we
would have \(\deg^-(y)\le1\), a contradiction; so \(y\notin N(x)\),
\(y\ne x\), hence \(y\) is the antipode \(\bar x\) of \(x\) (in \(Q_3\)
the only vertex at distance 3). Then \(N(y)\) is the set of the 3
vertices at distance 2 from \(x\), all lower than \(y\), so
\(\deg^-(y)=3\) and by (b) every \(b\in N(y)\) has \(\deg^-(b)\le1\).

\(N(x)\cup N(y)\) is the set \(S\) of the 6 vertices other than \(x,y\).
Count edges inside \(S\): \(Q_3\) has 12 edges, 3 are incident to \(x\),
3 to \(y\), and \(xy\) is not an edge, so exactly \(12-6=6\) edges have
both ends in \(S\). Each such edge is oriented towards its higher
endpoint, which lies in \(S\), so \(\sum_{s\in S}\deg^-(s)\ge 6\). But
each \(s\in S\) has \(\deg^-(s)\le1\), so
\(\sum_{s\in S}\deg^-(s)\le6\); therefore \(\deg^-(s)=1\) for every
\(s\in S\). Thus no vertex of \(S\) is a valley; \(x\) and \(y\) are not
valleys either (\(\deg^-\ge2\)). So the labelling has no valley,
contradicting the fact that the vertex with label 1 is a valley (all its
neighbours have larger labels). \(\square\)

Combining: \(U(Q_3)=14\), attained by the labelling above.

\subparagraph{Independent confirmation: exhaustive exact
computation}\label{independent-confirmation-exhaustive-exact-computation}

Script \texttt{q3\_esaustivo.py} (saved in \texttt{runs/p2\_q3/sandbox/}
and copied to \texttt{problema-2/certificati/}) enumerates \textbf{all}
\(8!=40320\) bijections \(V(Q_3)\to\{1,\dots,8\}\)
(\texttt{itertools.permutations}), with vertices encoded as integers
\(0..7\) (bit \(i\) = coordinate \(i\)) and adjacency = XOR with a
single bit. For each labelling it computes the number of uphill paths in
two independent ways --- (i) an explicit recursive DFS that enumerates
every uphill path from every valley, (ii) the DP of Lemma 1 processed in
increasing label order --- and \texttt{assert}s that the two agree. All
arithmetic is Python integers (exact; rigor = \texttt{exact}).

Output (Python 3.14.7, wall-clock 0.36 s):

\begin{verbatim}
etichettature esaminate: 40320
U(Q_3) = 14; etichettature ottime: 7104
distribuzione (cammini -> numero etichettature): [(14, 7104), (16, 14112), (17, 6720), (18, 4704), (19, 576), (20, 3648), (21, 1920), (22, 480), (23, 192), (24, 768), (26, 96)]
prima etichettatura ottima (ordine crescente di etichetta): ['000', '100', '010', '110', '101', '011', '001', '111']
\end{verbatim}

The distribution confirms Proposition 4 (no labelling with 13, and also
none with 15) and Corollary 3 (none below 13). Reproduce with
\texttt{.venv/bin/python\ problema-2/certificati/q3\_esaustivo.py}. A
second script \texttt{verifica\_etichettatura\_q3.py} checks only the
submitted labelling (bijection check, valleys, per-vertex \(p\), total
14) and prints the table above.

\subparagraph{\texorpdfstring{2b. Dimostrazione per
\(Q_4\)}{2b. Dimostrazione per Q\_4}}\label{b.-dimostrazione-per-q_4}

\paragraph{\texorpdfstring{\(U(Q_4)=34\)}{U(Q\_4)=34}}\label{uq_434}

\textbf{Conventions.} \(V(Q_4)=\{0,1\}^4\), \(u\sim v\) iff they differ
in exactly one coordinate; \(|V|=16\), \(|E|=32\), every vertex has
degree 4. In a 0/1 string the \(i\)-th character (from the left) is
coordinate \(i\); the convention is irrelevant for the count, since
every coordinate permutation is an automorphism of \(Q_4\) and
automorphisms preserve the number of uphill paths (they map valleys to
valleys and uphill paths to uphill paths bijectively when the labelling
is transported). For a labelling \(f\) orient each edge from the smaller
to the larger label (\(u\to v\) iff \(uv\in E\), \(f(u)<f(v)\));
\(\deg^-(v)=\#\{u:u\to v\}\), \(\deg^+(v)=\#\{w: v\to w\}\), so
\(\deg^-(v)+\deg^+(v)=4\). Let \(p(v)\) be the number of uphill paths
whose last vertex is \(v\).

\subparagraph{Preliminaries (verified claims, restated with proof for
self-containedness)}\label{preliminaries-verified-claims-restated-with-proof-for-self-containedness}

\textbf{Lemma 1.} \(p(v)=[v\text{ is a valley}]+\sum_{u\to v}p(u)\), and
the total number of uphill paths is \(\sum_v p(v)\).

\emph{Proof.} An uphill path ending at \(v\) either has length \(1\) ---
then it is \((v)\) and \(v\) is a valley, and conversely a valley \(v\)
gives exactly this path --- or has length \(\ge2\); deleting \(v\) gives
an uphill path ending at its penultimate vertex \(u\), with \(u\to v\),
and conversely appending \(v\) to an uphill path ending at \(u\) with
\(u\to v\) gives an uphill path ending at \(v\) (the increasing
condition is exactly \(f(u)<f(v)\)). These maps are mutually inverse,
and different \(u\) give disjoint sets (they differ in the penultimate
vertex). Each uphill path has a unique last vertex, hence the total is
\(\sum_v p(v)\). \(\square\)

\textbf{Lemma 2.} For every vertex \(v\): \(p(v)\ge1\) and
\(p(v)\ge\deg^-(v)\).

\emph{Proof.} From \(v\) repeatedly move to a neighbour with a smaller
label while one exists; the labels strictly decrease, so the walk stops,
at a vertex with no smaller neighbour, i.e.~a valley \(v_1\). The
reversed walk is an uphill path ending at \(v\), so \(p(v)\ge1\).
Applying this to every \(u\) with \(u\to v\) and using Lemma 1,
\(p(v)\ge\sum_{u\to v}p(u)\ge\sum_{u\to v}1=\deg^-(v)\). \(\square\)

\textbf{Corollary 3.} Total
\(\ \ge\sum_v\max(1,\deg^-(v))=\#\{\text{valleys}\}+|E|\).

\emph{Proof.} Termwise from Lemma 2 and Lemma 1. Valleys are exactly the
vertices with \(\deg^-=0\), and \(\sum_v\deg^-(v)=|E|\) because every
edge has exactly one head. So
\(\sum_v\max(1,\deg^-(v))=\#\{v:\deg^-(v)=0\}+\sum_{v:\deg^-(v)\ge1}\deg^-(v)=\#\{\text{valleys}\}+|E|\).
\(\square\)

The vertex with label \(1\) is always a valley (all its neighbours have
larger labels). Hence for \(Q_4\): total \(\ge 32+1=33\).

\subparagraph{\texorpdfstring{Lower bound: no labelling of \(Q_4\) has
exactly 33 uphill paths, hence
\(U(Q_4)\ge34\)}{Lower bound: no labelling of Q\_4 has exactly 33 uphill paths, hence U(Q\_4)\textbackslash ge34}}\label{lower-bound-no-labelling-of-q_4-has-exactly-33-uphill-paths-hence-uq_4ge34}

Suppose some labelling \(f\) of \(Q_4\) has total exactly \(33\).

\textbf{Step 1 (all inequalities are tight).} By Corollary 3,
\(33=\text{total}\ge\#\{\text{valleys}\}+32\ge 33\), so
\(\#\{\text{valleys}\}=1\) and \(\sum_v p(v)=\sum_v\max(1,\deg^-(v))\).
Since \(p(v)\ge\max(1,\deg^-(v))\) holds for every \(v\) (Lemma 2),
equality of the sums forces equality for every term:
\(p(v)=\max(1,\deg^-(v))\) for all \(v\). In particular, for every
non-valley \(v\) (i.e.~\(\deg^-(v)\ge1\)): \(p(v)=\deg^-(v)\).

\textbf{Step 2 (every in-neighbour of a non-valley has \(p=1\), hence
\(\deg^-\le1\)).} Let \(v\) be a non-valley. By Lemma 1 (with
\([v\text{ valley}]=0\)) and Step 1,
\[\sum_{u\to v}p(u)=p(v)=\deg^-(v)=\sum_{u\to v}1 .\] Each term
satisfies \(p(u)\ge1\) (Lemma 2), and there are \(\deg^-(v)\) terms on
each side, so \(p(u)=1\) for every \(u\to v\). By Lemma 2 again,
\(\deg^-(u)\le p(u)=1\).

\textbf{Step 3 (in-degrees lie in \(\{0,1,4\}\)).} Let \(u\) be any
vertex. If \(\deg^+(u)\ge1\), pick \(w\) with \(u\to w\); then \(w\) has
an in-neighbour, so \(w\) is not a valley, and Step 2 applied to \(v=w\)
gives \(\deg^-(u)\le1\). If \(\deg^+(u)=0\), then
\(\deg^-(u)=4-\deg^+(u)=4\). Hence \(\deg^-(u)\in\{0,1,4\}\) for every
vertex \(u\).

\textbf{Step 4 (counting edges).} Let \(s=\#\{u:\deg^-(u)=4\}\). The
vertices with \(\deg^-(u)=0\) are exactly the valleys, and there is
exactly one (Step 1). The remaining \(16-1-s=15-s\) vertices have
\(\deg^-(u)=1\). Summing in-degrees over all vertices, and using
\(\sum_u\deg^-(u)=|E|=32\):
\[4s+1\cdot(15-s)+0\cdot 1=32\quad\Longrightarrow\quad 3s=17,\] which
has no integer solution. Contradiction.

Therefore no labelling of \(Q_4\) has exactly \(33\) uphill paths;
combined with total \(\ge33\) (Corollary 3), every labelling has at
least \(34\) uphill paths: \(U(Q_4)\ge34\).

(\emph{Remark, not needed for the cell:} the same argument for \(Q_d\),
\(d\ge2\), gives total \(=|E|+1\) only if \((d-1)s=2^{d-1}(d-2)+1\) has
an integer solution; for \(d=3\) this is \(2s=5\), recovering
Proposition 4 of the \(Q_3\) proof.)

\subparagraph{Upper bound: an explicit labelling with exactly 34 uphill
paths}\label{upper-bound-an-explicit-labelling-with-exactly-34-uphill-paths}

Labelling \(f\), as the list of the 16 vertices in increasing label
order (label \(1\) first):

\[\texttt{0000},\ \texttt{1111},\ \texttt{0111},\ \texttt{1101},\ \texttt{1110},\ \texttt{1011},\ \texttt{0001},\ \texttt{1000},\ \texttt{0010},\ \texttt{0100},\ \texttt{1001},\ \texttt{0110},\ \texttt{0011},\ \texttt{1010},\ \texttt{0101},\ \texttt{1100}.\]

Structure: label \(1\) = weight \(0\); label \(2\) = weight \(4\);
labels \(3\)--\(6\) = the four vertices of weight \(3\); labels
\(7\)--\(10\) = the four vertices of weight \(1\); labels \(11\)--\(16\)
= the six vertices of weight \(2\). (This is a bijection: the sixteen
strings are pairwise distinct and are all of \(\{0,1\}^4\).)

Computation of \(p\) in increasing label order (Lemma 1). Neighbours of
a vertex of weight \(k\) have weights \(k\pm1\).

{\def\LTcaptype{none} % do not increment counter
\begin{longtable}[]{@{}
  >{\raggedright\arraybackslash}p{(\linewidth - 6\tabcolsep) * \real{0.2500}}
  >{\raggedright\arraybackslash}p{(\linewidth - 6\tabcolsep) * \real{0.2500}}
  >{\raggedright\arraybackslash}p{(\linewidth - 6\tabcolsep) * \real{0.2500}}
  >{\raggedright\arraybackslash}p{(\linewidth - 6\tabcolsep) * \real{0.2500}}@{}}
\toprule\noalign{}
\begin{minipage}[b]{\linewidth}\raggedright
vertex
\end{minipage} & \begin{minipage}[b]{\linewidth}\raggedright
label
\end{minipage} & \begin{minipage}[b]{\linewidth}\raggedright
lower neighbours
\end{minipage} & \begin{minipage}[b]{\linewidth}\raggedright
\(p\)
\end{minipage} \\
\midrule\noalign{}
\endhead
\bottomrule\noalign{}
\endlastfoot
0000 & 1 & none (its neighbours are the weight-1 vertices, labels 7--10)
--- \textbf{valley} & 1 \\
1111 & 2 & none (its neighbours are the weight-3 vertices, labels 3--6)
--- \textbf{valley} & 1 \\
0111 & 3 & 1111 (label 2); the other neighbours 0011, 0101, 0110 have
labels 13, 15, 12 & 1 \\
1101 & 4 & 1111; others 0101, 1001, 1100 have labels 15, 11, 16 & 1 \\
1110 & 5 & 1111; others 0110, 1010, 1100 have labels 12, 14, 16 & 1 \\
1011 & 6 & 1111; others 0011, 1001, 1010 have labels 13, 11, 14 & 1 \\
0001 & 7 & 0000 (label 1); others 1001, 0101, 0011 have labels 11, 15,
13 & 1 \\
1000 & 8 & 0000; others 1100, 1010, 1001 have labels 16, 14, 11 & 1 \\
0010 & 9 & 0000; others 1010, 0110, 0011 have labels 14, 12, 13 & 1 \\
0100 & 10 & 0000; others 1100, 0110, 0101 have labels 16, 12, 15 & 1 \\
1001 & 11 & 0001, 1000 (weight 1) and 1101, 1011 (weight 3): all four
lower, each with \(p=1\) & 4 \\
0110 & 12 & 0010, 0100, 1110, 0111 & 4 \\
0011 & 13 & 0001, 0010, 1011, 0111 & 4 \\
1010 & 14 & 1000, 0010, 1110, 1011 & 4 \\
0101 & 15 & 0001, 0100, 1101, 0111 & 4 \\
1100 & 16 & 1000, 0100, 1110, 1101 & 4 \\
\end{longtable}
}

Every weight-2 vertex has exactly two weight-1 and two weight-3
neighbours, all with labels \(\le10<11\), so all four neighbours are
lower and each has \(p=1\), giving \(p=4\). The valleys are exactly 0000
and 1111 (every other vertex has a lower neighbour, as the table shows).
Total \[\sum_v p(v)=1+1+4\cdot1+4\cdot1+6\cdot4=34 .\] Hence
\(U(Q_4)\le34\).

\subparagraph{Conclusion}\label{conclusion}

\[U(Q_4)=34,\] attained by the labelling above. (For comparison
\(U(Q_3)=14=|E|+2\) and \(U(Q_4)=34=|E|+2\).)

\subparagraph{Computational checks (exact integer arithmetic;
independent of the
proof)}\label{computational-checks-exact-integer-arithmetic-independent-of-the-proof}

\begin{enumerate}
\def\labelenumi{\arabic{enumi}.}
\item
  \texttt{verifica\_etichettatura\_q4.py} checks the submitted
  labelling: bijection, valleys, per-vertex \(p\) (the table above), and
  the total by two independent methods --- the DP of Lemma 1 in label
  order and an explicit recursive DFS enumerating every uphill path from
  every valley. Output: both totals \(=34\), valleys \(\{0000,1111\}\).
  Wall-clock 0.02 s (Python 3.14.7, macOS).
\item
  \texttt{q4\_branch\_and\_bound.py\ T} is an exact, symmetry-reduced
  search for labellings with \(<T\) uphill paths. Vertices are placed in
  increasing label order; when \(v\) is placed all its lower neighbours
  are already placed, so \(p(v)\) is fixed at that moment (Lemma 1).
  \emph{Symmetry reduction:} at each node only one candidate per orbit
  of the pointwise stabiliser (inside \(\mathrm{Aut}(Q_4)\), order
  \(384\), generated explicitly as coordinate permutation followed by
  XOR with a mask) of the already-placed set is tried; this is sound
  because if \(g\) fixes the placed vertices pointwise and \(g(v)=v'\),
  then \(f\mapsto f\circ g^{-1}\) maps labellings with prefix
  \((\dots,v)\) bijectively onto labellings with prefix \((\dots,v')\)
  and preserves the number of uphill paths. \emph{Lower bound at a node}
  (placed set \(S\), unplaced set \(U\),
  \(a(v)=\sum_{u\in N(v)\cap S}p(u)\), \(d(v)\) = number of
  in-neighbours of \(v\) inside \(U\) in the final orientation, so
  \(\sum_{v\in U}d(v)=E(U)\), the number of edges inside \(U\)): for
  \(v\in U\) all \(S\)-neighbours are lower and all unplaced
  in-neighbours have \(p\ge1\), so \(p(v)\ge\max(1,a(v)+d(v))\);
  minimising \(\sum_{v\in U}\max(1,a(v)+d(v))\) over nonnegative \(d\)
  with \(\sum d=E(U)\) gives
  \(\sum_{v\in U,a(v)\ge1}a(v)+\max(\#\{v\in U:a(v)=0\},E(U))\), so
  \(\text{LB}=\sum_{v\in S}p(v)+\sum_{v\in U,a(v)\ge1}a(v)+\max(\#\{a=0\},E(U))\)
  is a valid lower bound for every completion; a node is pruned iff
  \(\text{LB}\ge T\). Results: \(T=34\): \textbf{0 labellings found}, 47
  978 nodes, 0.3 s. \(T=35\): 19 525 440 labellings with exactly 34
  reached by the reduced search, 54 814 700 nodes, 156.8 s (shows the
  pruning is not over-aggressive at the optimum). Sanity check on
  \(Q_3\) (same code with \texttt{D\ =\ 3}): \(T=13\): 0 found;
  \(T=14\): 0 found (83 nodes); \(T=15\): 148 found, and
  \(148\cdot48=7104\) is exactly the known number of optimal labellings
  of \(Q_3\) (\(|\mathrm{Aut}(Q_3)|=48\)). This search is a confirmation
  only; the proof of \(U(Q_4)\ge34\) is the hand argument above.
\item
  \texttt{q4\_ricerca\_locale.py} (simulated annealing, 12 seeds, 200
  000 steps each, 12.8 s) found the value 34 in every run and never
  below; heuristic only, proves nothing, used to find the incumbent.
\end{enumerate}

\subparagraph{3. Verifica: istruzioni, dipendenze,
tempi}\label{verifica-istruzioni-dipendenze-tempi}

Tre programmi indipendenti, solo Python 3 standard (interi esatti),
ciascuno enumera tutte le \(8!=40320\) etichettature: -
\texttt{problema-2/certificati/q3\_esaustivo.py} --- DFS esplicita dei
cammini + programmazione dinamica di Lemma 1, con \texttt{assert} di
uguaglianza su ogni etichettatura. Tempo misurato: 0.37 s. -
\texttt{problema-2/certificati/q3\_verifica\_indipendente.py} ---
scritto separatamente, metodo diverso: costruzione esplicita livello per
livello di tutte le sequenze crescenti a partire dalle valli. Tempo
misurato: 0.20 s. -
\texttt{problema-2/certificati/verifica\_etichettatura\_q3.py} ---
controlla la sola etichettatura consegnata (biiezione, valli, \(p(v)\)
per vertice, totale 14). Comando:
\texttt{python3\ problema-2/certificati/q3\_verifica\_indipendente.py}
(e analoghi). Python 3.14.7, macOS. Tutti riportano: minimo 14, 7104
etichettature ottime, distribuzione senza valori 13 e 15.

Per \(Q_4\) (la prova del lower bound è a mano; il calcolo è solo
conferma): - \texttt{verifica\_etichettatura\_q4.py} (Researcher, due
metodi) e \texttt{q4\_verifica\_etichettatura\_indipendente.py} (scritto
a parte, costruzione esplicita dei cammini): entrambi 34 cammini, valli
\texttt{0000} e \texttt{1111}, biiezione verificata. -
\texttt{q4\_branch\_and\_bound.py\ SOGLIA}: branch-and-bound esatto con
riduzione per simmetria (384 automorfismi); con soglia 34 nessuna
etichettatura sotto 34 (47978 nodi, 0.3 s); con soglia 35 trova
etichettature con 34 (54.8M nodi, 157 s). La correttezza della riduzione
è argomentata nello script ma la cella non ne dipende. -
\texttt{q4\_ricerca\_locale.py}: simulated annealing, solo esplorazione
(12 seed, miglior valore 34 in ogni run).

\subparagraph{4. Fonti e contributo}\label{fonti-e-contributo}

Nessuna fonte esterna usata. L'argomento del lower bound generalizza
l'idea di IMO 2022 P6 (ogni vertice è fine di almeno un cammino;
\(p(v)\ge\deg^-(v)\)), aggiungendo l'ostruzione specifica di \(Q_3\)
(Proposizione 4). Prova prodotta dal Researcher automatico (Claude Fable
5.1, run \texttt{runs/p2\_q3/attempts/attempt\_001.json}) e revisionata
a mano passo per passo; verifica numerica indipendente scritta a parte.

\subparagraph{5. Limiti e osservazioni}\label{limiti-e-osservazioni}

\begin{itemize}
\tightlist
\item
  Nessuna parte irrisolta per C1.
\item
  \(U(Q_3)=|E|+2\) e \(U(Q_4)=|E|+2\). L'argomento di parità di \(Q_4\)
  (Step 3--4) si generalizza: il totale \(|E|+1\) è impossibile su
  \(Q_d\) quando \((d-1)\nmid 2^{d-1}(d-2)+1\), cioè per ogni \(d\) da 3
  a 10 almeno. Quindi \(U(Q_d)\ge|E|+2\) in quel range: bound debole
  rispetto ai 2368 noti per \(Q_9\), ma gratuito.
\item
  L'ottimo di \(Q_4\) ha DUE valli e realizza esattamente
  \(|E|+\#\text{valli}\): per i cubi grandi conviene cercare
  etichettature con poche valli in cui quasi tutti i vertici abbiano
  \(p(v)=\max(1,\deg^-(v))\).
\end{itemize}

\subsubsection{Decisioni degli agenti}
\paragraph{Run \texttt{p2\_c1}}
\paragraph{Run \texttt{p2\_q3}}
\begin{description}
\item[Iterazione 1 — attempt\_001]
\textbf{Researcher}: famiglia \texttt{exact\_computation}, sotto-obiettivo: Determine U(Q\_3) exactly (value + explicit optimal labelling in the required format), with certainty via exhaustive exact enumeration of all 8! labellings, plus an independent hand proof of the lower bound.. Stato dichiarato: \texttt{CELL\_SOLVED\_CANDIDATE}.
\emph{Perché questa scelta}: No prior attempts exist. The cell explicitly allows exhaustive exact computation with code and stated finite set; 8! = 40320 labellings is tiny (0.36 s). To make the result robust against a coding bug, two genuinely different counting routines are cross-checked on every labelling, and the lower bound is also proved by hand so the value does not rest on the code alone.
\textbf{Referee}: verdetto \texttt{(nessun verdetto)}.
\end{description}
\paragraph{Run \texttt{p2\_q4}}
\begin{description}
\item[Iterazione 1 — attempt\_001]
\textbf{Researcher}: famiglia \texttt{direct\_proof}, sotto-obiettivo: Determine U(Q\_4) exactly: prove the lower bound U(Q\_4) \textgreater{}= 34 by hand and exhibit an explicit labelling with exactly 34 uphill paths (route (b) of the cell), with an independent exact computer confirmation.. Stato dichiarato: \texttt{CELL\_SOLVED\_CANDIDATE}.
\emph{Perché questa scelta}: No failed attempts recorded. Simulated annealing (12 seeds) never went below 34 = |E| + 2, while the verified bound gives 33; the gap of 1 suggested a tightness analysis of the verified bound |E| + \#valleys, exactly as Proposition 4 did for Q\_3. The tightness analysis turned out to close with a pure counting argument (no case analysis), so I chose route (b) (hand proof + explicit labelling), and added a symmetry-reduced exact search only as an independent check, not as the proof.
\textbf{Referee}: verdetto \texttt{(nessun verdetto)}.
\end{description}

\subsubsection{Codice su cui poggia il tentativo}
Nessun codice nei tentativi.
\begin{lstlisting}[caption={attempt\_001 code\_1 [python, rigore exact] — Exhaustive exact enumeration of all 8! = 40320 labellings of Q\_3; two independent counts ( — verdetto del Referee: (nessun verdetto)}]
"""
Calcolo esatto di U(Q_3) per enumerazione esaustiva di tutte le 8! = 40320 etichettature.

Due conteggi indipendenti per ogni etichettatura:
  1. DFS: enumera esplicitamente ogni cammino in salita (v_1,...,v_k) partendo da ogni valle.
  2. DP:  p(v) = [v valle] + somma_{u vicino, f(u)<f(v)} p(u);  totale = somma_v p(v).
Aritmetica: solo interi Python (esatta). Insieme coperto: TUTTE le biiezioni V(Q_3) -> {1..8}.
"""
from itertools import permutations
from collections import Counter
import time

D = 3
VERTICI = list(range(2 ** D))                      # vertice = intero, bit i = coordinata i
VICINI = [[v ^ (1 << i) for i in range(D)] for v in VERTICI]


def stringa(v):
    """Vertice intero -> stringa 0/1 di lunghezza D (bit 0 a sinistra)."""
    return "".join(str((v >> i) & 1) for i in range(D))


def valli(etichetta):
    """Vertici con tutti i vicini di etichetta maggiore."""
    return [v for v in VERTICI if all(etichetta[w] > etichetta[v] for w in VICINI[v])]


def conta_dfs(etichetta):
    """Conteggio 1: enumera esplicitamente ogni cammino in salita."""
    def cammini_da(v):
        # il cammino che termina in v conta 1, piu' tutte le estensioni ai vicini maggiori
        return 1 + sum(cammini_da(w) for w in VICINI[v] if etichetta[w] > etichetta[v])
    return sum(cammini_da(v) for v in valli(etichetta))


def conta_dp(etichetta):
    """Conteggio 2: programmazione dinamica in ordine crescente di etichetta."""
    ordine = sorted(VERTICI, key=lambda v: etichetta[v])
    p = {}
    for v in ordine:
        minori = [w for w in VICINI[v] if etichetta[w] < etichetta[v]]
        p[v] = (1 if not minori else 0) + sum(p[w] for w in minori)
    return sum(p.values())


def esplora_tutte():
    """Enumera tutte le 8! etichettature; restituisce (minimo, distribuzione, ottime)."""
    distribuzione = Counter()
    ottime = []
    minimo = None
    for perm in permutations(range(1, 2 ** D + 1)):
        etichetta = list(perm)                     # etichetta[v] = f(v)
        n_dfs = conta_dfs(etichetta)
        n_dp = conta_dp(etichetta)
        assert n_dfs == n_dp, (etichetta, n_dfs, n_dp)
        distribuzione[n_dfs] += 1
        if minimo is None or n_dfs < minimo:
            minimo, ottime = n_dfs, [etichetta]
        elif n_dfs == minimo:
            ottime.append(etichetta)
    return minimo, distribuzione, ottime


def lista_ordinata(etichetta):
    """Formato di consegna: vertici in ordine crescente di etichetta, come stringhe 0/1."""
    return [stringa(v) for v in sorted(VERTICI, key=lambda v: etichetta[v])]


if __name__ == "__main__":
    inizio = time.time()
    minimo, distribuzione, ottime = esplora_tutte()
    print(f"etichettature esaminate: {sum(distribuzione.values())}")
    print(f"U(Q_3) = {minimo}; etichettature ottime: {len(ottime)}")
    print("distribuzione (cammini -> numero etichettature):", sorted(distribuzione.items()))
    print("prima etichettatura ottima (ordine crescente di etichetta):", lista_ordinata(ottime[0]))
    print(f"tempo: {time.time() - inizio:.2f} s")

\end{lstlisting}
\begin{lstlisting}[caption={attempt\_001 code\_2 [python, rigore exact] — Independent check of the single submitted labelling (bijection check, valleys, per-vertex  — verdetto del Referee: (nessun verdetto)}]
"""
Verifica indipendente della singola etichettatura consegnata per Q_3:
stampa valli, p(v) per ogni vertice e il totale dei cammini in salita (interi esatti).
"""
D = 3
CONSEGNA = ['000', '100', '010', '110', '101', '011', '001', '111']   # ordine crescente di etichetta


def vicini(s):
    """Stringhe 0/1 che differiscono da s in esattamente una coordinata."""
    return [s[:i] + ('1' if s[i] == '0' else '0') + s[i + 1:] for i in range(D)]


def verifica(consegna):
    """Calcola p(v) in ordine di etichetta e restituisce (valli, p, totale)."""
    assert sorted(consegna) == sorted(format(k, f'0{D}b') for k in range(2 ** D)), "non e' una biiezione"
    etichetta = {s: k + 1 for k, s in enumerate(consegna)}
    p, valli = {}, []
    for s in consegna:
        minori = [w for w in vicini(s) if etichetta[w] < etichetta[s]]
        if not minori:
            valli.append(s)
        p[s] = (1 if not minori else 0) + sum(p[w] for w in minori)
    return valli, p, sum(p.values())


if __name__ == "__main__":
    valli, p, totale = verifica(CONSEGNA)
    print("valli:", valli)
    for s in CONSEGNA:
        print(f"  {s}  etichetta {CONSEGNA.index(s)+1}  p = {p[s]}")
    print("cammini in salita totali:", totale)

\end{lstlisting}
\begin{lstlisting}[caption={attempt\_001 code\_1 [python, rigore exact] — Exact verification of the submitted labelling of Q\_4: bijection check, valleys, per-vertex — verdetto del Referee: (nessun verdetto)}]
"""Verifica ESATTA (interi Python) dell'etichettatura candidata di Q_4 con 34 cammini in salita.
Due metodi indipendenti: (i) DP di Lemma 1 in ordine di etichetta; (ii) enumerazione esplicita
(DFS) di tutti i cammini in salita da ogni valle. Stampa la tabella p(v) per la consegna."""
D = 4
N = 1 << D
# Convenzione: carattere i-esimo della stringa (da sinistra) = coordinata i; bit i dell'intero.
ETICHETTATURA = ['0000', '1111', '0111', '1101', '1110', '1011', '0001', '1000',
                 '0010', '0100', '1001', '0110', '0011', '1010', '0101', '1100']


def da_stringa(s):
    """Converte la stringa 0/1 (carattere i = coordinata i) nell'intero con bit i = coordinata i."""
    return sum(int(c) << i for i, c in enumerate(s))


def vicini(v):
    """I d vicini di v in Q_d: si cambia esattamente una coordinata."""
    return [v ^ (1 << i) for i in range(D)]


def conta_con_dp(ordine, etichetta):
    """Metodo (i): p(v) = [v valle] + somma p(u) sui vicini u con etichetta minore."""
    p = {}
    righe = []
    for v in ordine:
        minori = [u for u in vicini(v) if etichetta[u] < etichetta[v]]
        p[v] = 1 if not minori else sum(p[u] for u in minori)
        righe.append((format(v, f'0{D}b')[::-1], etichetta[v], [format(u, f'0{D}b')[::-1] for u in minori], p[v]))
    return sum(p.values()), righe


def conta_con_dfs(etichetta):
    """Metodo (ii): enumera esplicitamente ogni cammino in salita partendo da ogni valle."""
    def estendi(v):
        totale = 1  # il cammino che termina qui
        for w in vicini(v):
            if etichetta[w] > etichetta[v]:
                totale += estendi(w)
        return totale
    valli = [v for v in range(N) if all(etichetta[w] > etichetta[v] for w in vicini(v))]
    return sum(estendi(v) for v in valli), valli


def main():
    ordine = [da_stringa(s) for s in ETICHETTATURA]
    assert sorted(ordine) == list(range(N)), "non e' una biiezione"
    etichetta = {v: i + 1 for i, v in enumerate(ordine)}
    totale_dp, righe = conta_con_dp(ordine, etichetta)
    totale_dfs, valli = conta_con_dfs(etichetta)
    for stringa, lab, minori, p in righe:
        print(f"{stringa}  etichetta {lab:2d}  vicini minori {minori}  p = {p}")
    print("valli:", [format(v, f'0{D}b')[::-1] for v in valli])
    print("totale (DP di Lemma 1):", totale_dp)
    print("totale (DFS esplicita):", totale_dfs)
    assert totale_dp == totale_dfs == 34


if __name__ == "__main__":
    main()

\end{lstlisting}
\begin{lstlisting}[caption={attempt\_001 code\_2 [python, rigore exact] — Independent exact confirmation of the lower bound: symmetry-reduced branch-and-bound over  — verdetto del Referee: (nessun verdetto)}]
"""Branch-and-bound ESATTO (interi) su Q_4: cerca un'etichettatura con < SOGLIA cammini in salita.
Si piazzano i vertici in ordine crescente di etichetta; p(v) e' determinato al piazzamento (Lemma 1).
Riduzione per simmetria: a ogni nodo si prova un solo candidato per orbita dello stabilizzatore
puntuale (in Aut(Q_4), ordine 384) dell'insieme gia' piazzato. Bound inferiore dimostrato nel testo.
Conferma indipendente della prova a mano; NON e' la prova principale."""
import sys, time
from itertools import permutations

D = 4
N = 1 << D
VICINI = [[v ^ (1 << i) for i in range(D)] for v in range(N)]


def automorfismi():
    """Tutti i 384 automorfismi di Q_4 come tuple immagine: v -> perm(coordinate) XOR maschera."""
    gruppo = []
    for perm in permutations(range(D)):
        for maschera in range(N):
            immagine = tuple(sum(((v >> i) & 1) << perm[i] for i in range(D)) ^ maschera for v in range(N))
            gruppo.append(immagine)
    return gruppo


GRUPPO = automorfismi()


def bound_inferiore(p, piazzati, non_piazzati):
    """Bound: per v non piazzato p(v) >= max(1, a(v) + d(v)), a(v) = somma p sui vicini piazzati,
    d(v) = archi entranti da vertici non piazzati (somma dei d(v) = archi interni E(U)).
    Minimizzando sui d(v) a somma fissata: somma_{a>=1} a(v) + max(|U0|, E(U)), U0 = {a(v)=0}."""
    totale = sum(p[v] for v in piazzati)
    zeri, archi_interni = 0, 0
    for v in non_piazzati:
        a = sum(p[u] for u in VICINI[v] if u in piazzati)
        if a == 0:
            zeri += 1
        else:
            totale += a
        archi_interni += sum(1 for u in VICINI[v] if u not in piazzati)
    return totale + max(zeri, archi_interni // 2)


def rappresentanti(non_piazzati, stabilizzatore):
    """Un candidato per orbita dello stabilizzatore sui vertici non piazzati."""
    visti, scelti = set(), []
    for v in sorted(non_piazzati):
        if v in visti:
            continue
        scelti.append(v)
        visti.update(g[v] for g in stabilizzatore)
    return scelti


class Ricerca:
    def __init__(self, soglia):
        self.soglia = soglia
        self.nodi = 0
        self.trovate = []

    def esplora(self, ordine, p, piazzati, stabilizzatore):
        """Espande un nodo: prova ogni rappresentante d'orbita come prossimo vertice."""
        self.nodi += 1
        non_piazzati = [v for v in range(N) if v not in piazzati]
        if not non_piazzati:
            self.trovate.append(list(ordine))
            return
        if bound_inferiore(p, piazzati, non_piazzati) >= self.soglia:
            return
        for v in rappresentanti(non_piazzati, stabilizzatore):
            minori = [u for u in VICINI[v] if u in piazzati]
            p[v] = 1 if not minori else sum(p[u] for u in minori)
            piazzati.add(v)
            ordine.append(v)
            self.esplora(ordine, p, piazzati, [g for g in stabilizzatore if g[v] == v])
            ordine.pop()
            piazzati.discard(v)
            del p[v]


def main():
    soglia = int(sys.argv[1]) if len(sys.argv) > 1 else 34
    inizio = time.time()
    ricerca = Ricerca(soglia)
    ricerca.esplora([], {}, set(), GRUPPO)
    print(f"soglia {soglia}: etichettature con < {soglia} cammini trovate: {len(ricerca.trovate)}")
    for ordine in ricerca.trovate[:5]:
        print("  ", [format(v, f'0{D}b')[::-1] for v in ordine])
    print(f"nodi visitati: {ricerca.nodi}; tempo: {time.time() - inizio:.1f} s")


if __name__ == "__main__":
    main()

\end{lstlisting}
\begin{lstlisting}[caption={attempt\_001 code\_3 [python, rigore float\_exploration\_only] — Heuristic simulated annealing to find an incumbent labelling (upper bound candidate). 12 s — verdetto del Referee: (nessun verdetto)}]
"""Ricerca locale (simulated annealing) di etichettature di Q_4 con pochi cammini in salita.
Serve solo a trovare un buon candidato (upper bound): NON dimostra nulla.
Aritmetica: interi esatti (il conteggio), float solo nella regola di accettazione di SA."""
import random, sys, math

D = 4
N = 1 << D
VICINI = [[v ^ (1 << i) for i in range(D)] for v in range(N)]


def conta_cammini(ordine):
    """Numero di cammini in salita dell'etichettatura data come lista di vertici
    in ordine crescente di etichetta (Lemma 1: p(v) = [valle] + somma p(u) sui vicini minori)."""
    etichetta = [0] * N
    for pos, v in enumerate(ordine):
        etichetta[v] = pos
    p = [0] * N
    totale = 0
    for v in ordine:
        minori = [u for u in VICINI[v] if etichetta[u] < etichetta[v]]
        p[v] = 1 if not minori else sum(p[u] for u in minori)
        totale += p[v]
    return totale


def ricottura(seme, passi=200000):
    """Una corsa di simulated annealing con mosse di scambio e di spostamento."""
    rng = random.Random(seme)
    ordine = list(range(N))
    rng.shuffle(ordine)
    valore = conta_cammini(ordine)
    migliore, migliore_ordine = valore, ordine[:]
    for passo in range(passi):
        temperatura = 2.0 * (1 - passo / passi) + 0.05
        nuovo = ordine[:]
        i, j = rng.randrange(N), rng.randrange(N)
        if rng.random() < 0.5:
            nuovo[i], nuovo[j] = nuovo[j], nuovo[i]
        else:
            nuovo.insert(j, nuovo.pop(i))
        nuovo_valore = conta_cammini(nuovo)
        if nuovo_valore <= valore or rng.random() < math.exp((valore - nuovo_valore) / temperatura):
            ordine, valore = nuovo, nuovo_valore
            if valore < migliore:
                migliore, migliore_ordine = valore, ordine[:]
    return migliore, migliore_ordine


def main():
    corse = int(sys.argv[1]) if len(sys.argv) > 1 else 20
    globale, globale_ordine = None, None
    for seme in range(corse):
        valore, ordine = ricottura(seme)
        if globale is None or valore < globale:
            globale, globale_ordine = valore, ordine
        print(f"seme {seme}: {valore} (migliore finora {globale})", flush=True)
    stringhe = [format(v, f"0{D}b") for v in globale_ordine]
    print("MIGLIORE:", globale, stringhe)


if __name__ == "__main__":
    main()

\end{lstlisting}

\subsection{Cella 2 (2 punti) — stato: bozza pronta}
\subparagraph{\texorpdfstring{Parte 2 (C2) ---
\(U(Q_5)\)}{Parte 2 (C2) --- U(Q\_5)}}\label{parte-2-c2-uq_5}

\textbf{Punteggio:} 2 points · \textbf{Valutazione:} Checked instantly

\(Q_5\) has \(32\) vertices and \(80\) edges. Determine \(U(Q_5)\).

\textbf{Enunciato protetto dato agli agenti (state.json).} \texttt{Determine U(Q\_5) exactly and exhibit an optimal labelling of Q\_5 in the required format.}

\textbf{Evidenza registrata in STATUS.md.} \$U(Q\_5)=88\$: riduzione a (★) riletta a mano + 3 enumerazioni concordi; etichettatura verificata; Referee automatico: giudice A PASS (2 run), giudice B PARTIAL solo per regole segnaposto → terzo run con regole ufficiali in corso; bozza in submission/parte-2.md

\subsubsection{Consegna (prova)}
\paragraph{Problema 2 --- Parte 2 (C2) --- bozza di
consegna}\label{problema-2-parte-2-c2-bozza-di-consegna}

\textbf{Stato dichiarato: BOZZA PRONTA, da revisionare.} Valore
\(U(Q_5)=88\) con etichettatura esplicita; lower bound per riduzione a
un enunciato finito (★) confutato per enumerazione esaustiva (tre
implementazioni concordi, una indipendente per metodo). Verdetto del
Referee automatico: giudice A (matematica) PASS su entrambe le direzioni
in due run indipendenti; giudice B (evidenze): vedi §5. Revisione umana
della riduzione: fatta lemma per lemma (STATUS.md).

\subparagraph{1. Risultato e ambito}\label{risultato-e-ambito}

\(U(Q_5) = 88 = |E(Q_5)| + 8\). Etichettatura ottima (ordine crescente
di etichetta; posizione \(i\) della stringa = coordinata \(i\);
convenzione irrilevante perché le permutazioni di coordinate sono
automorfismi di \(Q_5\)), con 8 valli:
\texttt{00000,\ 11110,\ 11001,\ 10101,\ 01101,\ 10011,\ 01011,\ 00111}
sono le etichette 1--8 (le valli), seguite dalle foglie e dai picchi
come nella lista completa in §2 (tabella dell'upper bound). La lista
completa a 32 elementi è prodotta e verificata da
\texttt{certificati/costruzione\_88.py} e
\texttt{certificati/verifica\_dfs\_88.py}.

\subparagraph{2. Dimostrazione}\label{dimostrazione}

\emph{(Testo del tentativo
\texttt{runs/p2\_q5/attempts/attempt\_001.json} del Researcher, riletto
lemma per lemma dal team; le verifiche computazionali sono state
rieseguite dall'orchestratore in
\texttt{runs/p2\_q5/verifica/attempt\_001/}.)}

\paragraph{\texorpdfstring{\(U(Q_5) = 88\)}{U(Q\_5) = 88}}\label{uq_5-88}

\textbf{Conventions.} \(V(Q_5)=\{0,1\}^5\), \(u\sim v\) iff they differ
in exactly one coordinate; \(|V|=32\), \(|E|=80\), every vertex has
degree \(5\). In a 0/1 string the \(i\)-th character from the left is
coordinate \(i\) (irrelevant for the count: coordinate permutations are
automorphisms). For a labelling \(f\), orient each edge from the smaller
to the larger label (\(u\to v\) iff \(uv\in E\), \(f(u)<f(v)\));
\(\deg^-(v)=\#\{u:u\to v\}\), \(\deg^+(v)=\#\{w:v\to w\}\),
\(\deg^-(v)+\deg^+(v)=5\). Let \(p(v)\) be the number of uphill paths
whose last vertex is \(v\). A \emph{valley} is a vertex with
\(\deg^-=0\); a \emph{peak} is a vertex with \(\deg^+=0\). The vertex
with label \(1\) is a valley, so \(v:=\#\text{valleys}\ge1\).

\subparagraph{0. Preliminaries (restated with proof for
self-containedness)}\label{preliminaries-restated-with-proof-for-self-containedness}

\textbf{Lemma 1 (recursion).}
\(p(v)=[v\text{ valley}]+\sum_{u\to v}p(u)\) and the total number of
uphill paths is \(T=\sum_v p(v)\).

\emph{Proof.} An uphill path ending at \(v\) has length \(1\) --- then
it is \((v)\) and \(v\) is a valley; conversely a valley gives exactly
this path --- or length \(\ge2\); deleting \(v\) gives an uphill path
ending at its penultimate vertex \(u\) with \(u\to v\); conversely
appending \(v\) to an uphill path ending at \(u\) with \(u\to v\) gives
an uphill path ending at \(v\) (the increasing condition is exactly
\(f(u)<f(v)\)). These maps are mutually inverse and different \(u\) give
disjoint sets. Every uphill path has a unique last vertex. \(\square\)

\textbf{Lemma 2.} For every vertex \(v\): \(p(v)\ge1\), and
\(p(v)\ge\deg^-(v)\).

\emph{Proof.} From \(v\) move repeatedly to a neighbour with smaller
label while one exists; labels strictly decrease, so this stops at a
valley; the reversed walk is an uphill path ending at \(v\). Hence
\(p(v)\ge1\) for all \(v\), and by Lemma 1,
\(p(v)\ge\sum_{u\to v}p(u)\ge\deg^-(v)\). \(\square\)

\subparagraph{1. The excess identity}\label{the-excess-identity}

Define the \textbf{excess} \(X:=\sum_{(u\to w)\in E}\bigl(p(u)-1\bigr)\)
(sum over all \(80\) oriented edges). By Lemma 2 every term is \(\ge0\).

\textbf{Lemma 3.} \(T=|E|+v+X = 80+v+X\).

\emph{Proof.} By Lemma 1,
\(T=\sum_v p(v)=v+\sum_{w}\sum_{u\to w}p(u)=v+\sum_{(u\to w)\in E}\bigl(1+(p(u)-1)\bigr)=v+|E|+X\),
since each edge is counted once, at its head. \(\square\)

Grouping the terms of \(X\) by tail: \(X=\sum_u \deg^+(u)\,(p(u)-1)\).
Call \(u\) \textbf{heavy} if \(\deg^+(u)\ge1\) and \(p(u)\ge2\); let
\(H\) be the set of heavy vertices and \(c(u):=\deg^+(u)(p(u)-1)\) its
contribution, so \(X=\sum_{u\in H}c(u)\). Vertices with \(p=1\) are
called \textbf{light}.

\textbf{Lemma 4 (structure of light vertices).} If \(p(u)=1\) then
either \(u\) is a valley or \(\deg^-(u)=1\) and its unique in-neighbour
is light. Consequently, for any set \(L\) of light vertices, the
subgraph of \(Q_5\) induced by \(L\) is acyclic.

\emph{Proof.} By Lemma 2, \(\deg^-(u)\le p(u)=1\). If \(\deg^-(u)=1\)
with in-neighbour \(u'\), Lemma 1 gives \(1=p(u)=p(u')\). For the second
claim: a cycle inside \(L\) contains a vertex of maximal label, which
has two in-neighbours on the cycle, contradicting \(\deg^-\le1\).
\(\square\)

\textbf{Lemma 5 (every heavy vertex costs at least 3).} For every heavy
\(u\), \(c(u)\ge3\). Moreover, writing \(i=\deg^-(u)\) (so
\(\deg^+(u)=5-i\), \(1\le i\le4\) since \(\deg^+\ge1\) and \(p\ge2\)
excludes valleys): - if \(i=1\): \(p(u)=p(u')\ge2\) for the in-neighbour
\(u'\), which is then itself heavy (\(\deg^+(u')\ge1\) because
\(u'\to u\)), and \(c(u)=4(p(u)-1)\ge4\); - if \(i=2\):
\(c(u)=3(p(u)-1)\ge3\), with equality iff \(p(u)=2\) iff both
in-neighbours are light; - if \(i=3\): \(c(u)=2(p(u)-1)\ge4\) (as
\(p(u)\ge3\)), with equality iff all three in-neighbours are light; - if
\(i=4\): \(c(u)=p(u)-1\ge3\), with equality iff all four in-neighbours
are light.

\emph{Proof.} Direct from Lemma 2 (\(p(u)\ge i\)) and Lemma 1
(\(p(u)=\sum_{u'\to u}p(u')\) with each \(p(u')\ge1\); equality
\(p(u)=i\) iff every in-neighbour has \(p=1\)). Note also: if an
in-neighbour \(u'\) of a heavy \(u\) has \(p(u')\ge2\) then \(u'\) is
heavy (it has \(\deg^+\ge1\)). \(\square\)

\textbf{Corollary 6.} \(X\in\{0\}\cup[3,\infty)\), and \(X\le 7\) forces
\(|H|\le2\).

\textbf{Lemma 7 (in-degree count).} Let \(P\) be the set of peaks,
\(q=|P|\), \(s_H=\sum_{u\in H}\deg^-(u)\). If every non-heavy non-peak
vertex is light (which holds automatically for every vertex with
\(\deg^+\ge1\) that is not heavy, by definition of heavy), then
\[4q=48+v+|H|-s_H .\]

\emph{Proof.} Peaks have \(\deg^-=5\); no peak is heavy (heavy needs
\(\deg^+\ge1\)) and no valley is heavy. The vertices outside \(P\cup H\)
are light: valleys (\(\deg^-=0\), \(v\) of them) and light non-valleys
(\(\deg^-=1\) by Lemma 4, \(32-q-|H|-v\) of them). Summing in-degrees,
\(80=\sum_u\deg^-(u)=5q+s_H+(32-q-|H|-v)\), i.e.~\(4q=48+v+|H|-s_H\).
\(\square\)

Also note: \textbf{peaks form an independent set} (if two peaks were
adjacent, the one with the smaller label would have \(\deg^+\ge1\)).

\subparagraph{\texorpdfstring{2. Case analysis: every labelling with
\(T\le87\) has one of four
structures}{2. Case analysis: every labelling with T\textbackslash le87 has one of four structures}}\label{case-analysis-every-labelling-with-tle87-has-one-of-four-structures}

Assume \(T\le87\), i.e.~(Lemma 3) \(v+X\le7\) with \(v\ge1\), so
\(X\le6\). By Corollary 6, \(X\in\{0,3,4,5,6\}\) and \(|H|\le2\).

\textbf{Case \(|H|=2\).} Then \(X\ge6\), so \(X=6\), \(v=1\), and both
heavy vertices have \(c=3\); by Lemma 5 each has
\((\deg^-,p)\in\{(2,2),(4,4)\}\), so \(s_H\in\{4,6,8\}\) and Lemma 7
gives \(4q=48+1+2-s_H\in\{47,45,43\}\), none divisible by \(4\).
\textbf{Impossible.}

\textbf{Case \(|H|=0\).} Then \(X=0\), \(T=80+v\), and Lemma 7 gives
\(4q=48+v\), so \(v\equiv0\pmod 4\); with \(T\le87\) this forces
\(v=4\), \(q=13\), \(T=84\). Every non-peak vertex is light; so
\(F:=V\setminus P\) (19 vertices) induces an acyclic subgraph (Lemma 4),
and \(P\) is an independent set of size \(13\). \textbf{(Structure A.)}

\textbf{Case \(|H|=1\), \(H=\{u\}\).} Then \(X=c(u)\le6\). By Lemma 5,
if some in-neighbour of \(u\) had \(p\ge2\) it would be heavy,
contradicting \(|H|=1\); hence all in-neighbours of \(u\) are light and
\(p(u)=\deg^-(u)=:i\); \(i=1\) is excluded (it needs a heavy
in-neighbour). So \((i,c(u))\in\{(2,3),(3,4),(4,3)\}\), and Lemma 7 with
\(|H|=1\), \(s_H=i\) gives \(4q=49+v-i\): - \(i=2\): \(4q=47+v\),
\(v\equiv1\pmod4\), and \(v+X=v+3\le7\) gives \(v=1\), \(q=12\),
\(T=84\). \textbf{(Structure B.)} - \(i=3\): \(4q=46+v\),
\(v\equiv2\pmod4\), \(v+4\le7\) gives \(v=2\), \(q=12\), \(T=86\).
\textbf{(Structure C.)} - \(i=4\): \(4q=45+v\), \(v\equiv3\pmod4\),
\(v+3\le7\) gives \(v=3\), \(q=12\), \(T=86\). \textbf{(Structure D.)}

In each of B, C, D: the out-neighbours \(w\) of \(u\) satisfy
\(p(w)\ge p(u)\ge2\) (Lemma 1), so if such a \(w\) had \(\deg^+(w)\ge1\)
it would be heavy; hence \textbf{all out-neighbours of \(u\) are peaks}.
The in-neighbours of \(u\) are non-peaks. All vertices other than \(u\)
and the peaks are light, so by Lemma 4 the set
\(L:=V\setminus(P\cup\{u\})\) induces an acyclic subgraph; \(P\) is an
independent set of size \(12\) and \(u\notin P\).

\textbf{Summary.} If \(T\le 87\) then there exist an independent set
\(I\subseteq V(Q_5)\) and a vertex \(u\notin I\) such that
\(V\setminus(I\cup\{u\})\) induces an acyclic subgraph, where either
\(|I|=12\) (structures B, C, D with \(I=P\)), or \(|I|=13\) and
\(V\setminus I\) is acyclic (structure A) --- in which case, choosing
any \(u_0\in I\) and \(I'=I\setminus\{u_0\}\), the pair \((I',u_0)\) has
\(|I'|=12\), \(I'\) independent, \(u_0\notin I'\), and
\(V\setminus(I'\cup\{u_0\})=V\setminus I\) acyclic. So in all cases:

\begin{quote}
\textbf{(★)} there is an independent set \(I\) of \(Q_5\) with
\(|I|=12\) and a vertex \(u\notin I\) such that \(Q_5-(I\cup\{u\})\) is
a forest.
\end{quote}

\subparagraph{\texorpdfstring{3. Exact finite verification: (★) never
holds in
\(Q_5\)}{3. Exact finite verification: (★) never holds in Q\_5}}\label{exact-finite-verification-never-holds-in-q_5}

\textbf{Reduction.} The maps \(x\mapsto x\oplus t\) are automorphisms of
\(Q_5\) and act transitively on vertices, and (★) is invariant under
automorphisms. Hence it suffices to check all independent \(12\)-sets
\(I\) containing the vertex \(00000\) and all \(u\notin I\).

\textbf{Computation 1} (\texttt{verifica\_q5\_indipendenti.py}, exact
integer arithmetic): backtracking enumerates all independent sets of
size \(12\) containing \(0\) in increasing vertex order (a vertex may be
added iff it is neither in the set nor adjacent to it), and for each of
the \(20\) vertices \(u\notin I\) tests acyclicity of the induced
subgraph on the \(19\) remaining vertices by the exact criterion
\emph{\#edges = \#vertices − \#components}. Output:
\texttt{insiemi\ indipendenti\ di\ taglia\ 12\ contenenti\ 0:\ 1377;\ coppie\ (I,u)\ esaminate:\ 27540;\ complementi\ aciclici\ trovati:\ 0;\ tempo\ 0.24s}.

\textbf{Computation 2} (\texttt{verifica\_q5\_bipartita.py}, independent
method, no symmetry reduction): every independent set is \(A\cup B\)
with \(A\) a subset of the 16 even-weight vertices and \(B\) a subset of
the odd-weight vertices not adjacent to \(A\); the script scans all
\(2^{16}\) bitmasks \(A\), all \(B\) of size \(12-|A|\), and tests
acyclicity by union--find over the edge list. Output:
\texttt{insiemi\ indipendenti\ di\ taglia\ 12\ (tutti):\ 3672;\ coppie\ (I,u)\ esaminate:\ 73440;\ complementi\ aciclici:\ 0;\ tempo\ 0.39s}.
Consistency: \(3672\cdot 12/32=1377\), as vertex-transitivity predicts.

Hence (★) is false, so no labelling of \(Q_5\) has \(T\le87\):
\[U(Q_5)\ge 88.\]

(The counts also make the impossibility of the \(X=0\), \(v=4\)
structure transparent by hand: an independent \(12\)- or \(13\)-set
turns out to lie in one parity class except for a few mixed sets, and
any three even vertices contain two at distance \(2\), which share two
odd neighbours and thus close a \(4\)-cycle in the complement; but the
proof relies only on the exhaustive check above.)

\subparagraph{4. Upper bound: a labelling with exactly 88 uphill
paths}\label{upper-bound-a-labelling-with-exactly-88-uphill-paths}

Let \(R=\{00000,11110\}\) and let \(P\) be the \(14\) even-weight
vertices other than \(R\). \(P\) is independent (even-weight vertices
are pairwise non-adjacent). The complement
\(F=R\cup\{\text{16 odd vertices}\}\) induces exactly the \(10\) edges
from \(R\) to the odd vertices (the two vertices of \(R\) are at
distance \(4\), so their neighbourhoods
\(\{10000,01000,00100,00010,00001\}\) and
\(\{01110,10110,11010,11100,11111\}\) are disjoint): \(F\) is a forest
with \(18\) vertices, \(10\) edges and \(8\) components, namely the two
stars centred at \(R\) and the \(6\) isolated odd vertices
\(11001,10101,01101,10011,01011,00111\).

Labelling (increasing label order; coordinate 1 on the left):

\begin{verbatim}
00000,11110,11001,10101,01101,10011,01011,00111,
10000,01000,00100,11100,00010,11010,10110,01110,00001,11111,
11000,10100,01100,10010,01010,00110,10001,01001,00101,11101,00011,11011,10111,01111
\end{verbatim}

i.e.~labels \(1\)--\(2\): the star centres; \(3\)--\(8\): the six
isolated odd vertices; \(9\)--\(18\): the ten star leaves;
\(19\)--\(32\): the fourteen peaks \(P\).

\emph{Count.} The \(8\) vertices with labels \(1\)--\(8\) have all
neighbours higher (their \(F\)-neighbours, if any, are leaves labelled
\(9\)--\(18\); their other neighbours are in \(P\)): they are valleys,
\(p=1\). Each leaf (labels \(9\)--\(18\)) has exactly one lower
neighbour, its star centre, so \(p=1\). Each peak has all \(5\)
neighbours lower, all in \(F\) with \(p=1\), so \(p=5\). Total
\(T=18\cdot1+14\cdot5=88\). Verified exactly by two independent scripts:
\texttt{costruzione\_88.py} (recursion of Lemma 1) and
\texttt{verifica\_dfs\_88.py} (explicit DFS enumeration of all uphill
paths from the \(8\) valleys), both printing \(88\) and the same \(8\)
valleys.

\subparagraph{5. Conclusion}\label{conclusion}

\(U(Q_5)=88=|E|+8\), attained by the labelling in §4. Consistency checks
(not part of the proof): the same Lemmas 3--7 for \(Q_4\) give
\(c(u)\ge2\) and \(3q=16+v\), reproducing \(U(Q_4)=34\) with \(v=2\);
for \(Q_3\) they reproduce \(U(Q_3)=14\). The identity \(T=80+v+X\) and
the fact \(X\in\{0\}\cup[3,\infty)\) were also sanity-checked on
\(200{,}000\) random labellings of \(Q_5\) (\texttt{sanity\_eccesso.py};
the smallest positive \(X\) seen was \(14\)). Simulated annealing over
labellings (\texttt{ricottura\_q5.py}, 200k swap moves per seed, exact
scoring) never went below \(88\) in any seed run so far (evidence only).

\subparagraph{3. Verifica: istruzioni, dipendenze,
tempi}\label{verifica-istruzioni-dipendenze-tempi}

Dipendenze: Python 3 standard, nessun pacchetto. Tutti gli script usano
aritmetica intera esatta.

\begin{verbatim}
cd problema-2/certificati
python3 costruzione_88.py                 # costruisce l'etichettatura e conta 88 cammini (< 0.1 s)
python3 verifica_dfs_88.py                # conteggio indipendente per DFS esplicita (< 0.1 s)
python3 verifica_q5_indipendenti.py       # (★) confutata: 1377 insiemi con 0, 27540 coppie, 0 aciclici (0.47 s)
python3 verifica_q5_bipartita.py          # (★) confutata senza riduzione per simmetria: 3672 / 73440 / 0 (0.71 s)
python3 q5_stella_verifica_indipendente.py  # (★) confutata con metodo diverso (potatura delle foglie): 3672 / 73440 / 0 (0.40 s)
\end{verbatim}

Tempi misurati su portatile (riesecuzione fidata dell'orchestratore,
2026-09-26). Codice del tentativo: - \texttt{code\_1} (python, rigore
\texttt{exact}): Exact count of uphill paths of a labelling of Q\_d via
the recursion p(v) = {[}valley{]} + sum of p over lower neighbours
(shared module). - \texttt{code\_2} (python, rigore \texttt{exact}):
Certificate 1 for the lower bound: enumerate all independent 12-sets of
Q\_5 containing vertex 0 (1377 sets, symmetry reduction by XOR
translations) and, for eac - \texttt{code\_3} (python, rigore
\texttt{exact}): Certificate 2 for the lower bound, independent method
and no symmetry reduction: all independent 12-sets as A ∪ B (A ⊆ even
class over all 2\^{}16 bitmasks, B ⊆ od - \texttt{code\_4} (python,
rigore \texttt{exact}): Builds the 88-path labelling (peaks = even
vertices minus 00000, 11110) and counts exactly with the recursion:
prints 88, the 8 valleys and the 32-vertex order. - \texttt{code\_5}
(python, rigore \texttt{exact}): Independent verification of the
submitted labelling by explicit DFS enumeration of every uphill path
from each valley (string-based, no shared code): prints 88. -
\texttt{code\_6} (python, rigore \texttt{exact}): Sanity check (evidence
only, not part of the proof) of the identity T = 80 + \#valleys + X and
of X ∈ \{0\} ∪ {[}3,∞) on 200000 random labellings of Q\_5 (all asserts
- \texttt{code\_7} (python, rigore \texttt{float\_exploration\_only}):
Simulated annealing over labellings (exploration only): best value found
88 in every seed run (seeds 0-5 in a first run, seeds 0-3 in a second
run; 200k swap mo

\subparagraph{4. Fonti e contributo}\label{fonti-e-contributo}

\begin{itemize}
\tightlist
\item
  Lemma di conteggio \(p(v)=[\text{valle}]+\sum_{u\to v}p(u)\) e bound
  \(\#\ge|E|+\#\text{valli}\): nostri (parte 1, ridimostrati qui).
\item
  Contesto: IMO 2022 P6 (griglia), citato solo come contesto; nessun
  risultato preso dalla letteratura.
\item
  Ricerca arXiv deterministica (\texttt{runs/p2\_q5/literature.json},
  \texttt{problema-2/fonti/arxiv.json}): nessun lavoro sui cammini in
  salita sull'ipercubo; il problema risulta inedito. Contributo
  interamente nostro.
\end{itemize}

\subparagraph{5. Limiti e parti
irrisolte}\label{limiti-e-parti-irrisolte}

\begin{itemize}
\tightlist
\item
  Gap dichiarati dal Researcher nel tentativo:
\item
  The lower bound is computer-assisted at one point: statement (★) (no
  independent 12-set I and vertex u with Q\_5 − (I ∪ \{u\}) acyclic) is
  proved by exhaustive enumeration, not by hand. The reduction to (★) is
  fully by hand; the enumeration is exact, covers the whole finite set
  (all 3672 independent 12-sets × 20 outside vertices), runs in
  \textless{} 1 s and was done by two methods. A hand proof of (★) was
  only sketched (parity-class argument), not written out.
\item
  In Lemma 7 / §2 I use that every non-peak, non-heavy vertex is light;
  this is immediate from the definitions (a vertex with deg\^{}+ ≥ 1 and
  p ≥ 2 is heavy by definition) but the Referee should check that no
  case escapes: valleys (p=1), peaks (deg\^{}+=0), heavy, light ---
  these four classes cover V.
\item
  The step `out-neighbours of the unique heavy vertex u are peaks'
  (needed for nothing in the final finite check, since (★) only uses I
  independent of size 12 and u ∉ I) is included for completeness; the
  final reduction deliberately uses only the weaker necessary
  conditions, so any slip there does not affect the proof.
\item
  Simulated annealing results are exploration only and play no role in
  the proof; the second SA run (20 seeds) was still running when this
  report was written (4 seeds completed, all at 88).
\item
  Il giudice B ha classificato tutte le osservazioni su fonti e calcoli
  come INFO (fonti di contesto non caricate, versione di Python non
  registrata: 3.14, libreria standard); l'unico punto MISSING riguardava
  le regole ufficiali, ora passate alla lettera dal preambolo
  dell'enunciato.
\item
  La cella è ``checked instantly'' sul valore: il rischio residuo è un
  errore nella riduzione a (★); per questo la riduzione è stata riletta
  a mano e (★) confutata da tre programmi diversi.
\end{itemize}

\subsubsection{Decisioni degli agenti}
\paragraph{Run \texttt{p2\_c2}}
\paragraph{Run \texttt{p2\_q5}}
\begin{description}
\item[Iterazione 1 — attempt\_001]
\textbf{Researcher}: famiglia \texttt{case\_analysis}, sotto-obiettivo: . Stato dichiarato: \texttt{}.
\textbf{Referee}: verdetto \texttt{UNKNOWN\_STATUS} (review\_status \texttt{INCOMPLETE}).
\end{description}

\subsubsection{Codice su cui poggia il tentativo}
Nessun codice nei tentativi.
\begin{lstlisting}[caption={attempt\_001 code\_1 [python, rigore exact] — Exact count of uphill paths of a labelling of Q\_d via the recursion p(v) = [valley] + sum  — verdetto del Referee: UNKNOWN\_STATUS}]
"""Conteggio esatto dei cammini in salita di una etichettatura di Q_d.

Vertici = interi 0..2^d-1 (bit i = coordinata i); vicini = XOR con un bit.
p(v) = [v valle] + somma p(u) sui vicini u con etichetta minore; totale = somma p(v).
Aritmetica intera: risultato esatto.
"""


def vicini(v, d):
    """Restituisce i d vicini di v in Q_d."""
    return [v ^ (1 << i) for i in range(d)]


def conta_cammini(ordine, d):
    """ordine = lista dei vertici in ordine crescente di etichetta. Ritorna (totale, p, valli)."""
    etichetta = {v: i for i, v in enumerate(ordine)}
    assert sorted(ordine) == list(range(1 << d)), "non e' una biiezione"
    p = {}
    valli = []
    for v in ordine:
        minori = [u for u in vicini(v, d) if etichetta[u] < etichetta[v]]
        if not minori:
            valli.append(v)
            p[v] = 1
        else:
            p[v] = sum(p[u] for u in minori)
    return sum(p.values()), p, valli


def stringa(v, d):
    """Vertice intero -> stringa 0/1, coordinata 1 a sinistra."""
    return "".join(str((v >> i) & 1) for i in range(d))


def da_stringa(s):
    """Stringa 0/1 -> vertice intero (coordinata 1 a sinistra)."""
    return sum(int(c) << i for i, c in enumerate(s))

\end{lstlisting}
\noindent\emph{Riesecuzione fidata dell'orchestratore}: \texttt{orchestrator re-ran code\_1.py (python3, clean copy of the researcher sandbox): exit 0 in 0.0s; stdout: ''; stderr: ''}
\begin{lstlisting}[caption={attempt\_001 code\_2 [python, rigore exact] — Certificate 1 for the lower bound: enumerate all independent 12-sets of Q\_5 containing ver — verdetto del Referee: UNKNOWN\_STATUS}]
"""Certificato per il bound inferiore U(Q_5) >= 88 (implementazione 1: backtracking).

Enunciato verificato (Lemma computazionale):
  per ogni insieme indipendente I di Q_5 con |I| = 12 e per ogni vertice u di Q_5,
  il sottografo indotto da V meno (I u {u}) contiene un ciclo.
Riduzione: le traslazioni x -> x XOR t sono automorfismi di Q_5, quindi basta considerare
gli I che contengono il vertice 0 (ogni I non vuoto si trasla a contenerne uno).
Aciclicita' decisa con conteggio esatto: componenti + spigoli (un grafo e' una foresta sse
spigoli = vertici - componenti). Tutto intero: rigore = exact.
"""
import time
from conta_cammini import vicini

D = 5
N = 1 << D
ADJ = [set(vicini(v, D)) for v in range(N)]


def e_foresta(S):
    """True sse il sottografo indotto da S e' aciclico (spigoli = |S| - componenti)."""
    S = set(S)
    spigoli = sum(len(ADJ[v] & S) for v in S) // 2
    visti, comp = set(), 0
    for s in S:
        if s in visti:
            continue
        comp += 1
        pila, visti = [s], visti | {s}
        while pila:
            v = pila.pop()
            for w in ADJ[v] & S:
                if w not in visti:
                    visti.add(w)
                    pila.append(w)
    return spigoli == len(S) - comp


def indipendenti_con_zero(k):
    """Genera tutti gli insiemi indipendenti di taglia k contenenti 0 (vertici in ordine crescente)."""
    def ric(I, prossimo, vietati):
        if len(I) == k:
            yield list(I)
            return
        for v in range(prossimo, N):
            if v not in vietati and N - v >= k - len(I):
                yield from ric(I + [v], v + 1, vietati | ADJ[v] | {v})
    yield from ric([0], 1, ADJ[0] | {0})


if __name__ == "__main__":
    inizio = time.time()
    n_insiemi, n_test, violazioni = 0, 0, []
    for I in indipendenti_con_zero(12):
        n_insiemi += 1
        for u in range(N):
            if u in I:
                continue
            n_test += 1
            if e_foresta(set(range(N)) - set(I) - {u}):
                violazioni.append((I, u))
    print(f"insiemi indipendenti di taglia 12 contenenti 0: {n_insiemi}")
    print(f"coppie (I,u) esaminate: {n_test}; complementi aciclici trovati: {len(violazioni)}")
    print(f"tempo {time.time()-inizio:.2f}s")

\end{lstlisting}
\noindent\emph{Riesecuzione fidata dell'orchestratore}: \texttt{orchestrator re-ran code\_2.py (python3, clean copy of the researcher sandbox): exit 0 in 0.5s; stdout: 'insiemi indipendenti di taglia 12 contenenti 0: 1377\textbackslash{}ncoppie (I,u) esaminate: 27540; complementi aciclici trovati: 0\textbackslash{}ntempo 0.44s'; stderr: ''}
\begin{lstlisting}[caption={attempt\_001 code\_3 [python, rigore exact] — Certificate 2 for the lower bound, independent method and no symmetry reduction: all indep — verdetto del Referee: UNKNOWN\_STATUS}]
"""Certificato per il bound inferiore U(Q_5) >= 88 (implementazione 2, metodo diverso).

Stesso enunciato di verifica_q5_indipendenti.py, ma senza usare la simmetria:
ogni insieme indipendente e' A u B con A sottoinsieme dei vertici di peso pari e
B sottoinsieme dei vertici di peso dispari NON adiacenti ad A. Si scorrono tutti i 2^16
sottoinsiemi A (bitmask), poi tutti i B di taglia 12-|A| tra i dispari disponibili.
Aciclicita' con union-find (metodo diverso dal conteggio componenti). Rigore = exact.
"""
import itertools
import time

D = 5
N = 1 << D
PARI = [v for v in range(N) if bin(v).count("1") % 2 == 0]
DISPARI = [v for v in range(N) if bin(v).count("1") % 2 == 1]
SPIGOLI = [(v, v ^ (1 << i)) for v in range(N) for i in range(D) if v < v ^ (1 << i)]


def maschera_vicini(v):
    """Bitmask dei vicini di v."""
    return sum(1 << (v ^ (1 << i)) for i in range(D))


VIC = [maschera_vicini(v) for v in range(N)]


def ha_ciclo(maschera_vertici):
    """True sse il sottografo indotto dai vertici nella bitmask contiene un ciclo (union-find)."""
    padre = list(range(N))

    def trova(x):
        while padre[x] != x:
            padre[x] = padre[padre[x]]
            x = padre[x]
        return x

    for a, b in SPIGOLI:
        if (maschera_vertici >> a) & 1 and (maschera_vertici >> b) & 1:
            ra, rb = trova(a), trova(b)
            if ra == rb:
                return True
            padre[ra] = rb
    return False


def insiemi_indipendenti_12():
    """Genera le bitmask di tutti gli insiemi indipendenti di taglia 12."""
    for bits in range(1 << 16):
        A = [PARI[i] for i in range(16) if (bits >> i) & 1]
        if len(A) > 12:
            continue
        vietati = 0
        for a in A:
            vietati |= VIC[a]
        disponibili = [w for w in DISPARI if not (vietati >> w) & 1]
        for B in itertools.combinations(disponibili, 12 - len(A)):
            m = 0
            for x in A + list(B):
                m |= 1 << x
            yield m


if __name__ == "__main__":
    inizio = time.time()
    tutti = (1 << N) - 1
    n_insiemi, n_test, aciclici = 0, 0, 0
    for I in insiemi_indipendenti_12():
        n_insiemi += 1
        for u in range(N):
            if (I >> u) & 1:
                continue
            n_test += 1
            if not ha_ciclo(tutti & ~I & ~(1 << u)):
                aciclici += 1
    print(f"insiemi indipendenti di taglia 12 (tutti, senza simmetria): {n_insiemi}")
    print(f"coppie (I,u) esaminate: {n_test}; complementi aciclici: {aciclici}")
    print(f"tempo {time.time()-inizio:.2f}s")

\end{lstlisting}
\noindent\emph{Riesecuzione fidata dell'orchestratore}: \texttt{orchestrator re-ran code\_3.py (python3, clean copy of the researcher sandbox): exit 0 in 0.8s; stdout: 'insiemi indipendenti di taglia 12 (tutti, senza simmetria): 3672\textbackslash{}ncoppie (I,u) esaminate: 73440; complementi aciclici: 0\textbackslash{}ntempo 0.77s'; stderr: ''}
\begin{lstlisting}[caption={attempt\_001 code\_4 [python, rigore exact] — Builds the 88-path labelling (peaks = even vertices minus 00000, 11110) and counts exactly — verdetto del Referee: UNKNOWN\_STATUS}]
"""Costruzione esplicita con 88 cammini in salita su Q_5 e verifica esatta.

Idea: picchi = tutti i vertici di peso pari tranne 00000 e 11110 (14 picchi, indipendenti);
il complemento (2 pari + 16 dispari) e' una foresta con 10 spigoli e 8 componenti:
etichettiamo prima le 8 radici (valli), poi le foglie, poi i picchi. Ogni non-picco ha p = 1,
ogni picco ha p = 5: totale 18 + 14*5 = 88.
"""
from conta_cammini import conta_cammini, stringa, da_stringa

D = 5
N = 1 << D
peso = lambda v: bin(v).count("1")
radici = [da_stringa("00000"), da_stringa("11110")]
dispari = [v for v in range(N) if peso(v) % 2 == 1]
foglie = [v for v in dispari if any(v ^ r in [1 << i for i in range(D)] for r in radici)]
isolati = [v for v in dispari if v not in foglie]
picchi = [v for v in range(N) if peso(v) % 2 == 0 and v not in radici]
ordine = radici + isolati + foglie + picchi

if __name__ == "__main__":
    tot, p, valli = conta_cammini(ordine, D)
    print("totale cammini:", tot)
    print("valli:", [stringa(v, D) for v in valli])
    print("ordine:", ",".join(stringa(v, D) for v in ordine))

\end{lstlisting}
\noindent\emph{Riesecuzione fidata dell'orchestratore}: \texttt{orchestrator re-ran code\_4.py (python3, clean copy of the researcher sandbox): exit 0 in 0.0s; stdout: "totale cammini: 88\textbackslash{}nvalli: ['00000', '11110', '11001', '10101', '01101', '10011', '01011', '00111']\textbackslash{}nordine: 00000,11110,11001,10101,01101,10011,01011,00111,10000,01000,00100,11100,00010,11010,10110,01110,00001,11111,11000,10100,01100,10010,01010,00110,10001,01001,00101,11101,00011,11011,10111,0}
\begin{lstlisting}[caption={attempt\_001 code\_5 [python, rigore exact] — Independent verification of the submitted labelling by explicit DFS enumeration of every u — verdetto del Referee: UNKNOWN\_STATUS}]
"""Verifica indipendente dell'etichettatura a 88 cammini: enumerazione esplicita per DFS.

Metodo diverso dalla ricorsione p(v): da ogni valle si esplorano tutti i cammini con etichette
crescenti e si contano uno a uno. Controlla anche che l'input sia una biiezione.
"""
ORDINE = ("00000,11110,11001,10101,01101,10011,01011,00111,10000,01000,00100,11100,00010,11010,"
          "10110,01110,00001,11111,11000,10100,01100,10010,01010,00110,10001,01001,00101,11101,"
          "00011,11011,10111,01111").split(",")
D = 5


def vicini(s):
    """Vicini della stringa s: cambia un carattere."""
    return [s[:i] + ("1" if s[i] == "0" else "0") + s[i + 1:] for i in range(D)]


def conta_da(s, etichetta):
    """Numero di cammini in salita che iniziano in s (s incluso come cammino di lunghezza 1)."""
    return 1 + sum(conta_da(w, etichetta) for w in vicini(s) if etichetta[w] > etichetta[s])


if __name__ == "__main__":
    assert len(set(ORDINE)) == 32 and all(len(s) == 5 and set(s) <= {"0", "1"} for s in ORDINE)
    etichetta = {s: i + 1 for i, s in enumerate(ORDINE)}
    valli = [s for s in ORDINE if all(etichetta[w] > etichetta[s] for w in vicini(s))]
    totale = sum(conta_da(s, etichetta) for s in valli)
    print("valli:", valli)
    print("cammini in salita (DFS esplicita):", totale)

\end{lstlisting}
\noindent\emph{Riesecuzione fidata dell'orchestratore}: \texttt{orchestrator re-ran code\_5.py (python3, clean copy of the researcher sandbox): exit 0 in 0.0s; stdout: "valli: ['00000', '11110', '11001', '10101', '01101', '10011', '01011', '00111']\textbackslash{}ncammini in salita (DFS esplicita): 88"; stderr: ''}
\begin{lstlisting}[caption={attempt\_001 code\_6 [python, rigore exact] — Sanity check (evidence only, not part of the proof) of the identity T = 80 + \#valleys + X  — verdetto del Referee: UNKNOWN\_STATUS}]
"""Controllo di sanita' (solo evidenza) dell'identita' usata nella prova:
totale = |E| + #valli + X, X = somma sugli spigoli u->w di (p(u)-1), e X in {0} u [3, inf).
Etichettature casuali di Q_5 con aritmetica esatta.
"""
import random
from conta_cammini import conta_cammini, vicini

D, N = 5, 32
rng = random.Random(1)
valori_x = set()
for _ in range(200000):
    ordine = list(range(N))
    rng.shuffle(ordine)
    tot, p, valli = conta_cammini(ordine, D)
    et = {v: i for i, v in enumerate(ordine)}
    x = sum(p[u] - 1 for u in range(N) for w in vicini(u, D) if et[u] < et[w])
    assert tot == 80 + len(valli) + x
    valori_x.add(x)
print("identita' verificata su 200000 etichettature casuali; min X>0 osservato:", min(v for v in valori_x if v > 0))
print("valori di X osservati sotto 10:", sorted(v for v in valori_x if v < 10))

\end{lstlisting}
\noindent\emph{Riesecuzione fidata dell'orchestratore}: \texttt{orchestrator re-ran code\_6.py (python3, clean copy of the researcher sandbox): exit 0 in 13.4s; stdout: "identita' verificata su 200000 etichettature casuali; min X\textgreater{}0 osservato: 14\textbackslash{}nvalori di X osservati sotto 10: []"; stderr: ''}
\begin{lstlisting}[caption={attempt\_001 code\_7 [python, rigore float\_exploration\_only] — Simulated annealing over labellings (exploration only): best value found 88 in every seed  — verdetto del Referee: UNKNOWN\_STATUS}]
"""Ricerca locale (simulated annealing) sul numero di cammini in salita di Q_5.

Esplorazione: trova etichettature con pochi cammini. NON e' una prova (minimo locale possibile).
Stato = permutazione dei 32 vertici; mossa = scambio di due etichette; punteggio esatto (interi).
"""
import random
import sys
import time
from conta_cammini import conta_cammini, stringa

D = 5
N = 1 << D


def ricottura(seed, passi=200000, t0=2.0, t1=0.05):
    """Un run di annealing con seed fissato; ritorna (miglior_valore, miglior_ordine)."""
    rng = random.Random(seed)
    ordine = list(range(N))
    rng.shuffle(ordine)
    val = conta_cammini(ordine, D)[0]
    migliore, miglior_ordine = val, ordine[:]
    for k in range(passi):
        t = t0 * (t1 / t0) ** (k / passi)
        i, j = rng.randrange(N), rng.randrange(N)
        ordine[i], ordine[j] = ordine[j], ordine[i]
        nuovo = conta_cammini(ordine, D)[0]
        if nuovo <= val or rng.random() < 2.718 ** ((val - nuovo) / t):
            val = nuovo
            if val < migliore:
                migliore, miglior_ordine = val, ordine[:]
        else:
            ordine[i], ordine[j] = ordine[j], ordine[i]
    return migliore, miglior_ordine


if __name__ == "__main__":
    semi = int(sys.argv[1]) if len(sys.argv) > 1 else 8
    inizio = time.time()
    globale = None
    for seed in range(semi):
        v, o = ricottura(seed)
        print(f"seed {seed}: {v}", flush=True)
        if globale is None or v < globale[0]:
            globale = (v, o)
    v, o = globale
    tot, p, valli = conta_cammini(o, D)
    print("migliore:", v, "valli:", [stringa(x, D) for x in valli])
    print("ordine:", ",".join(stringa(x, D) for x in o))
    print(f"tempo {time.time()-inizio:.1f}s")

\end{lstlisting}
\noindent\emph{Riesecuzione fidata dell'orchestratore}: \texttt{orchestrator re-ran code\_7.py (python3, clean copy of the researcher sandbox): exit 0 in 60.5s; stdout: "seed 0: 88\textbackslash{}nseed 1: 88\textbackslash{}nseed 2: 88\textbackslash{}nseed 3: 88\textbackslash{}nseed 4: 88\textbackslash{}nseed 5: 88\textbackslash{}nseed 6: 88\textbackslash{}nseed 7: 88\textbackslash{}nmigliore: 88 valli: ['01001', '11100', '11010', '11111', '10011']\textbackslash{}nordine: 01001,11100,01000,11010,11111,01101,10100,10110,00110,00100,00111,01110,11001,01011,00001,11110,00010,10101,10011,10000,1100}

\subsection{Cella 3 (3 punti) — stato: in analisi}
\subparagraph{\texorpdfstring{Parte 3 (C3) ---
\(U(Q_6)\)}{Parte 3 (C3) --- U(Q\_6)}}\label{parte-3-c3-uq_6}

\textbf{Punteggio:} 3 points · \textbf{Valutazione:} Checked instantly

\(Q_6\) has \(64\) vertices and \(192\) edges. Determine \(U(Q_6)\).

\textbf{Enunciato protetto dato agli agenti (state.json).} \texttt{Determine U(Q\_6) exactly (64 vertices, 192 edges) and exhibit an optimal labelling in the required format. Checked on the value: it must be the true minimum, so a lower-bound proof (hand or certified computation \textless{} 10 min) is needed.}

\textbf{Evidenza registrata in STATUS.md.} —

\subsubsection{Consegna (prova)}
Nessuna bozza di consegna ancora scritta.

\subsubsection{Decisioni degli agenti}
\paragraph{Run \texttt{p2\_c3}}

\subsubsection{Codice su cui poggia il tentativo}
Nessun codice nei tentativi.

\subsection{Cella 4 (5 punti) — stato: in analisi}
\subparagraph{\texorpdfstring{Parte 4 (C4) --- \(U(Q_7)\) and
\(U(Q_8)\)}{Parte 4 (C4) --- U(Q\_7) and U(Q\_8)}}\label{parte-4-c4-uq_7-and-uq_8}

\textbf{Punteggio:} 5 points · \textbf{Valutazione:} Checked instantly

Two larger cubes: \(Q_7\) has \(128\) vertices and \(448\) edges,
\(Q_8\) has \(256\) vertices and \(1024\) edges. Determine \(U(Q_7)\)
and \(U(Q_8)\).

\textbf{Enunciato protetto dato agli agenti (state.json).} \texttt{Determine U(Q\_7) and U(Q\_8) exactly, each with an explicit optimal labelling in the required format. Checked on the values: lower-bound proofs needed.}

\textbf{Evidenza registrata in STATUS.md.} —

\subsubsection{Consegna (prova)}
Nessuna bozza di consegna ancora scritta.

\subsubsection{Decisioni degli agenti}
\paragraph{Run \texttt{p2\_c4}}

\subsubsection{Codice su cui poggia il tentativo}
Nessun codice nei tentativi.

\subsection{Cella 5 (8 punti) — stato: in analisi}
\subparagraph{\texorpdfstring{Parte 5 (C5) --- Bounds for
\(U(Q_9)\)}{Parte 5 (C5) --- Bounds for U(Q\_9)}}\label{parte-5-c5-bounds-for-uq_9}

\textbf{Punteggio:} 8 points · \textbf{Valutazione:} Judged

\(Q_9\) has \(512\) vertices and \(2304\) edges. The best bounds known
to the organisers are \[
2368 \le U(Q_9) \le 2400;
\] the lower bound is unpublished. Improve either one: prove that
\(U(Q_9) \ge 2369\), or exhibit a labelling of \(Q_9\) with at most
\(2399\) uphill paths.

\textbf{Enunciato protetto dato agli agenti (state.json).} \texttt{Improve one of the known bounds 2368 \textless{}= U(Q\_9) \textless{}= 2400: either prove U(Q\_9) \textgreater{}= 2369, or exhibit a labelling of Q\_9 with at most 2399 uphill paths (its paths are counted mechanically).}

\textbf{Evidenza registrata in STATUS.md.} via realistica: etichettatura \$\textbackslash{}le2399\$

\subsubsection{Consegna (prova)}
Nessuna bozza di consegna ancora scritta.

\subsubsection{Decisioni degli agenti}
\paragraph{Run \texttt{p2\_c5}}

\subsubsection{Codice su cui poggia il tentativo}
Nessun codice nei tentativi.

\subsection{Cella 6 (13 punti) — stato: in analisi}
\subparagraph{\texorpdfstring{Parte 6 (C6) ---
\(U(Q_9)\)}{Parte 6 (C6) --- U(Q\_9)}}\label{parte-6-c6-uq_9}

\textbf{Punteggio:} 13 points · \textbf{Valutazione:} Judged ·
\textbf{Open question}

The same cube, settled completely. Determine \(U(Q_9)\) exactly, with
proof of both bounds.

\emph{(Fine del problema 2. Punteggi: 1+2+3+5+8+13 = 32.)}

\textbf{Enunciato protetto dato agli agenti (state.json).} \texttt{OPEN. Determine U(Q\_9) exactly, with an optimal labelling and a proof of both bounds.}

\textbf{Evidenza registrata in STATUS.md.} —

\subsubsection{Consegna (prova)}
Nessuna bozza di consegna ancora scritta.

\subsubsection{Decisioni degli agenti}
\paragraph{Run \texttt{p2\_c6}}

\subsubsection{Codice su cui poggia il tentativo}
Nessun codice nei tentativi.

\subsection{Certificati di verifica indipendente del team}
\begin{lstlisting}[caption={certificati/costruzione\_88.py}]
"""Costruzione esplicita con 88 cammini in salita su Q_5 e verifica esatta.

Idea: picchi = tutti i vertici di peso pari tranne 00000 e 11110 (14 picchi, indipendenti);
il complemento (2 pari + 16 dispari) e' una foresta con 10 spigoli e 8 componenti:
etichettiamo prima le 8 radici (valli), poi le foglie, poi i picchi. Ogni non-picco ha p = 1,
ogni picco ha p = 5: totale 18 + 14*5 = 88.
"""
from conta_cammini import conta_cammini, stringa, da_stringa

D = 5
N = 1 << D
peso = lambda v: bin(v).count("1")
radici = [da_stringa("00000"), da_stringa("11110")]
dispari = [v for v in range(N) if peso(v) % 2 == 1]
foglie = [v for v in dispari if any(v ^ r in [1 << i for i in range(D)] for r in radici)]
isolati = [v for v in dispari if v not in foglie]
picchi = [v for v in range(N) if peso(v) % 2 == 0 and v not in radici]
ordine = radici + isolati + foglie + picchi

if __name__ == "__main__":
    tot, p, valli = conta_cammini(ordine, D)
    print("totale cammini:", tot)
    print("valli:", [stringa(v, D) for v in valli])
    print("ordine:", ",".join(stringa(v, D) for v in ordine))

\end{lstlisting}
\begin{lstlisting}[caption={certificati/q3\_esaustivo.py}]
"""
Calcolo esatto di U(Q_3) per enumerazione esaustiva di tutte le 8! = 40320 etichettature.

Due conteggi indipendenti per ogni etichettatura:
  1. DFS: enumera esplicitamente ogni cammino in salita (v_1,...,v_k) partendo da ogni valle.
  2. DP:  p(v) = [v valle] + somma_{u vicino, f(u)<f(v)} p(u);  totale = somma_v p(v).
Aritmetica: solo interi Python (esatta). Insieme coperto: TUTTE le biiezioni V(Q_3) -> {1..8}.
"""
from itertools import permutations
from collections import Counter
import time

D = 3
VERTICI = list(range(2 ** D))                      # vertice = intero, bit i = coordinata i
VICINI = [[v ^ (1 << i) for i in range(D)] for v in VERTICI]


def stringa(v):
    """Vertice intero -> stringa 0/1 di lunghezza D (bit 0 a sinistra)."""
    return "".join(str((v >> i) & 1) for i in range(D))


def valli(etichetta):
    """Vertici con tutti i vicini di etichetta maggiore."""
    return [v for v in VERTICI if all(etichetta[w] > etichetta[v] for w in VICINI[v])]


def conta_dfs(etichetta):
    """Conteggio 1: enumera esplicitamente ogni cammino in salita."""
    def cammini_da(v):
        # il cammino che termina in v conta 1, piu' tutte le estensioni ai vicini maggiori
        return 1 + sum(cammini_da(w) for w in VICINI[v] if etichetta[w] > etichetta[v])
    return sum(cammini_da(v) for v in valli(etichetta))


def conta_dp(etichetta):
    """Conteggio 2: programmazione dinamica in ordine crescente di etichetta."""
    ordine = sorted(VERTICI, key=lambda v: etichetta[v])
    p = {}
    for v in ordine:
        minori = [w for w in VICINI[v] if etichetta[w] < etichetta[v]]
        p[v] = (1 if not minori else 0) + sum(p[w] for w in minori)
    return sum(p.values())


def esplora_tutte():
    """Enumera tutte le 8! etichettature; restituisce (minimo, distribuzione, ottime)."""
    distribuzione = Counter()
    ottime = []
    minimo = None
    for perm in permutations(range(1, 2 ** D + 1)):
        etichetta = list(perm)                     # etichetta[v] = f(v)
        n_dfs = conta_dfs(etichetta)
        n_dp = conta_dp(etichetta)
        assert n_dfs == n_dp, (etichetta, n_dfs, n_dp)
        distribuzione[n_dfs] += 1
        if minimo is None or n_dfs < minimo:
            minimo, ottime = n_dfs, [etichetta]
        elif n_dfs == minimo:
            ottime.append(etichetta)
    return minimo, distribuzione, ottime


def lista_ordinata(etichetta):
    """Formato di consegna: vertici in ordine crescente di etichetta, come stringhe 0/1."""
    return [stringa(v) for v in sorted(VERTICI, key=lambda v: etichetta[v])]


if __name__ == "__main__":
    inizio = time.time()
    minimo, distribuzione, ottime = esplora_tutte()
    print(f"etichettature esaminate: {sum(distribuzione.values())}")
    print(f"U(Q_3) = {minimo}; etichettature ottime: {len(ottime)}")
    print("distribuzione (cammini -> numero etichettature):", sorted(distribuzione.items()))
    print("prima etichettatura ottima (ordine crescente di etichetta):", lista_ordinata(ottime[0]))
    print(f"tempo: {time.time() - inizio:.2f} s")

\end{lstlisting}
\begin{lstlisting}[caption={certificati/q3\_verifica\_indipendente.py}]
"""
Verifica indipendente di U(Q_3) = 14 (scritta separatamente dal Researcher, metodo diverso).

Metodo: per ogni etichettatura (8! = 40320) si costruisce l'insieme di tutte le sequenze crescenti di vertici
adiacenti per lunghezza crescente (BFS per livelli sull'ordine delle etichette), partendo dalle sole valli.
Nessuna ricorsione e nessuna programmazione dinamica: si contano le sequenze una per una, come oggetti espliciti.
Interi esatti. Uso: python3 problema-2/certificati/q3_verifica_indipendente.py
"""
import itertools
import time

VERTICI = list(range(8))                       # vertice = intero 0..7, bit i = coordinata i
VICINI = {v: [v ^ (1 << i) for i in range(3)] for v in VERTICI}


def valli(etichetta):
    """Vertici con tutti i vicini più alti: sono gli unici punti di partenza ammessi."""
    return [v for v in VERTICI if all(etichetta[w] > etichetta[v] for w in VICINI[v])]


def conta_cammini(etichetta):
    """Costruisce esplicitamente tutti i cammini in salita, livello per livello, e li conta."""
    livello = [(v,) for v in valli(etichetta)]  # cammini di lunghezza 1
    totale = 0
    while livello:
        totale += len(livello)
        livello = [c + (w,) for c in livello for w in VICINI[c[-1]] if etichetta[w] > etichetta[c[-1]]]
    return totale


def stringa(v):
    """Vertice come stringa 0/1 con la posizione i = coordinata i (bit 0 a sinistra)."""
    return "".join(str((v >> i) & 1) for i in range(3))


if __name__ == "__main__":
    inizio = time.time()
    minimo, esempio, quante_ottime = None, None, 0
    for perm in itertools.permutations(VERTICI):          # perm[k] = vertice con etichetta k+1
        etichetta = {v: k + 1 for k, v in enumerate(perm)}
        n = conta_cammini(etichetta)
        if minimo is None or n < minimo:
            minimo, esempio, quante_ottime = n, perm, 1
        elif n == minimo:
            quante_ottime += 1
    print(f"U(Q_3) = {minimo}; etichettature ottime: {quante_ottime}; tempo {time.time() - inizio:.2f} s")
    print("esempio ottimo (ordine crescente di etichetta):", [stringa(v) for v in esempio])
    # controllo della specifica etichettatura consegnata dal Researcher
    consegnata = ["000", "100", "010", "110", "101", "011", "001", "111"]
    et = {int(s[::-1], 2): k + 1 for k, s in enumerate(consegnata)}
    print("cammini dell'etichettatura consegnata:", conta_cammini(et))

\end{lstlisting}
\begin{lstlisting}[caption={certificati/q4\_branch\_and\_bound.py}]
"""Branch-and-bound ESATTO (interi) su Q_4: cerca un'etichettatura con < SOGLIA cammini in salita.
Si piazzano i vertici in ordine crescente di etichetta; p(v) e' determinato al piazzamento (Lemma 1).
Riduzione per simmetria: a ogni nodo si prova un solo candidato per orbita dello stabilizzatore
puntuale (in Aut(Q_4), ordine 384) dell'insieme gia' piazzato. Bound inferiore dimostrato nel testo.
Conferma indipendente della prova a mano; NON e' la prova principale."""
import sys, time
from itertools import permutations

D = 4
N = 1 << D
VICINI = [[v ^ (1 << i) for i in range(D)] for v in range(N)]


def automorfismi():
    """Tutti i 384 automorfismi di Q_4 come tuple immagine: v -> perm(coordinate) XOR maschera."""
    gruppo = []
    for perm in permutations(range(D)):
        for maschera in range(N):
            immagine = tuple(sum(((v >> i) & 1) << perm[i] for i in range(D)) ^ maschera for v in range(N))
            gruppo.append(immagine)
    return gruppo


GRUPPO = automorfismi()


def bound_inferiore(p, piazzati, non_piazzati):
    """Bound: per v non piazzato p(v) >= max(1, a(v) + d(v)), a(v) = somma p sui vicini piazzati,
    d(v) = archi entranti da vertici non piazzati (somma dei d(v) = archi interni E(U)).
    Minimizzando sui d(v) a somma fissata: somma_{a>=1} a(v) + max(|U0|, E(U)), U0 = {a(v)=0}."""
    totale = sum(p[v] for v in piazzati)
    zeri, archi_interni = 0, 0
    for v in non_piazzati:
        a = sum(p[u] for u in VICINI[v] if u in piazzati)
        if a == 0:
            zeri += 1
        else:
            totale += a
        archi_interni += sum(1 for u in VICINI[v] if u not in piazzati)
    return totale + max(zeri, archi_interni // 2)


def rappresentanti(non_piazzati, stabilizzatore):
    """Un candidato per orbita dello stabilizzatore sui vertici non piazzati."""
    visti, scelti = set(), []
    for v in sorted(non_piazzati):
        if v in visti:
            continue
        scelti.append(v)
        visti.update(g[v] for g in stabilizzatore)
    return scelti


class Ricerca:
    def __init__(self, soglia):
        self.soglia = soglia
        self.nodi = 0
        self.trovate = []

    def esplora(self, ordine, p, piazzati, stabilizzatore):
        """Espande un nodo: prova ogni rappresentante d'orbita come prossimo vertice."""
        self.nodi += 1
        non_piazzati = [v for v in range(N) if v not in piazzati]
        if not non_piazzati:
            self.trovate.append(list(ordine))
            return
        if bound_inferiore(p, piazzati, non_piazzati) >= self.soglia:
            return
        for v in rappresentanti(non_piazzati, stabilizzatore):
            minori = [u for u in VICINI[v] if u in piazzati]
            p[v] = 1 if not minori else sum(p[u] for u in minori)
            piazzati.add(v)
            ordine.append(v)
            self.esplora(ordine, p, piazzati, [g for g in stabilizzatore if g[v] == v])
            ordine.pop()
            piazzati.discard(v)
            del p[v]


def main():
    soglia = int(sys.argv[1]) if len(sys.argv) > 1 else 34
    inizio = time.time()
    ricerca = Ricerca(soglia)
    ricerca.esplora([], {}, set(), GRUPPO)
    print(f"soglia {soglia}: etichettature con < {soglia} cammini trovate: {len(ricerca.trovate)}")
    for ordine in ricerca.trovate[:5]:
        print("  ", [format(v, f'0{D}b')[::-1] for v in ordine])
    print(f"nodi visitati: {ricerca.nodi}; tempo: {time.time() - inizio:.1f} s")


if __name__ == "__main__":
    main()

\end{lstlisting}
\begin{lstlisting}[caption={certificati/q4\_ricerca\_locale.py}]
"""Ricerca locale (simulated annealing) di etichettature di Q_4 con pochi cammini in salita.
Serve solo a trovare un buon candidato (upper bound): NON dimostra nulla.
Aritmetica: interi esatti (il conteggio), float solo nella regola di accettazione di SA."""
import random, sys, math

D = 4
N = 1 << D
VICINI = [[v ^ (1 << i) for i in range(D)] for v in range(N)]


def conta_cammini(ordine):
    """Numero di cammini in salita dell'etichettatura data come lista di vertici
    in ordine crescente di etichetta (Lemma 1: p(v) = [valle] + somma p(u) sui vicini minori)."""
    etichetta = [0] * N
    for pos, v in enumerate(ordine):
        etichetta[v] = pos
    p = [0] * N
    totale = 0
    for v in ordine:
        minori = [u for u in VICINI[v] if etichetta[u] < etichetta[v]]
        p[v] = 1 if not minori else sum(p[u] for u in minori)
        totale += p[v]
    return totale


def ricottura(seme, passi=200000):
    """Una corsa di simulated annealing con mosse di scambio e di spostamento."""
    rng = random.Random(seme)
    ordine = list(range(N))
    rng.shuffle(ordine)
    valore = conta_cammini(ordine)
    migliore, migliore_ordine = valore, ordine[:]
    for passo in range(passi):
        temperatura = 2.0 * (1 - passo / passi) + 0.05
        nuovo = ordine[:]
        i, j = rng.randrange(N), rng.randrange(N)
        if rng.random() < 0.5:
            nuovo[i], nuovo[j] = nuovo[j], nuovo[i]
        else:
            nuovo.insert(j, nuovo.pop(i))
        nuovo_valore = conta_cammini(nuovo)
        if nuovo_valore <= valore or rng.random() < math.exp((valore - nuovo_valore) / temperatura):
            ordine, valore = nuovo, nuovo_valore
            if valore < migliore:
                migliore, migliore_ordine = valore, ordine[:]
    return migliore, migliore_ordine


def main():
    corse = int(sys.argv[1]) if len(sys.argv) > 1 else 20
    globale, globale_ordine = None, None
    for seme in range(corse):
        valore, ordine = ricottura(seme)
        if globale is None or valore < globale:
            globale, globale_ordine = valore, ordine
        print(f"seme {seme}: {valore} (migliore finora {globale})", flush=True)
    stringhe = [format(v, f"0{D}b") for v in globale_ordine]
    print("MIGLIORE:", globale, stringhe)


if __name__ == "__main__":
    main()

\end{lstlisting}
\begin{lstlisting}[caption={certificati/q4\_verifica\_etichettatura\_indipendente.py}]
"""
Verifica indipendente dell'etichettatura di Q_4 consegnata dal Researcher (metodo diverso dal suo):
costruzione esplicita, livello per livello, di tutti i cammini in salita a partire dalle valli.
Controlla anche che la lista sia una biiezione su {0,1}^4. Interi esatti.
Uso: python3 problema-2/certificati/q4_verifica_etichettatura_indipendente.py
"""
D = 4
CONSEGNATA = ["0000", "1111", "0111", "1101", "1110", "1011", "0001", "1000",
              "0010", "0100", "1001", "0110", "0011", "1010", "0101", "1100"]


def vicini(s):
    """Le d stringhe che differiscono da s in una sola posizione."""
    return [s[:i] + ("1" if s[i] == "0" else "0") + s[i + 1:] for i in range(D)]


def conta_cammini(etichetta):
    """Tutti i cammini in salita come tuple esplicite, per lunghezza crescente; ritorna il totale."""
    valli = [v for v in etichetta if all(etichetta[w] > etichetta[v] for w in vicini(v))]
    livello = [(v,) for v in valli]
    totale = 0
    while livello:
        totale += len(livello)
        livello = [c + (w,) for c in livello for w in vicini(c[-1]) if etichetta[w] > etichetta[c[-1]]]
    return totale, valli


if __name__ == "__main__":
    assert len(CONSEGNATA) == 16 and len(set(CONSEGNATA)) == 16 and all(len(s) == 4 and set(s) <= {"0", "1"} for s in CONSEGNATA)
    etichetta = {s: k + 1 for k, s in enumerate(CONSEGNATA)}   # etichetta = posizione nella lista
    totale, valli = conta_cammini(etichetta)
    print(f"biiezione ok; valli: {valli}; cammini in salita: {totale}")

\end{lstlisting}
\begin{lstlisting}[caption={certificati/q5\_stella\_verifica\_indipendente.py}]
"""
Verifica indipendente dell'affermazione (★) usata nel lower bound di U(Q_5) (metodo diverso dagli script del Researcher).

(★): esistono un insieme indipendente I di Q_5 con |I| = 12 e un vertice u ∉ I tali che Q_5 − (I ∪ {u}) è una foresta.
La prova mostra che ogni etichettatura con ≤ 87 cammini implica (★); qui si verifica che (★) è FALSA.

Metodo: enumerazione di tutti gli insiemi indipendenti di taglia 12 come bitmask su 32 bit (backtracking con
maschera dei vertici vietati), senza riduzione per simmetria; per ogni u ∉ I il test di aciclicità è per POTATURA
DELLE FOGLIE: si tolgono ripetutamente i vertici di grado ≤ 1 nel sottografo indotto; è una foresta se e solo se
non resta nulla. Interi esatti, nessuna dipendenza.  Uso: python3 problema-2/certificati/q5_stella_verifica_indipendente.py
"""
import time

D, N = 5, 32
ADJ = [sum(1 << (v ^ (1 << i)) for i in range(D)) for v in range(N)]   # vicini di v come bitmask
TUTTI = (1 << N) - 1


def insiemi_indipendenti(taglia):
    """Genera le bitmask di tutti gli insiemi indipendenti di `taglia` vertici (ordine crescente, senza ripetizioni)."""
    def ricorsione(insieme, vietati, prossimo, mancanti):
        if mancanti == 0:
            yield insieme
            return
        for v in range(prossimo, N - mancanti + 1):
            if not (vietati >> v) & 1:
                yield from ricorsione(insieme | (1 << v), vietati | ADJ[v] | (1 << v), v + 1, mancanti - 1)
    yield from ricorsione(0, 0, 0, taglia)


def e_foresta(vertici):
    """Potatura delle foglie sul sottografo indotto dalla bitmask `vertici`: resta vuoto ⇔ aciclico."""
    while vertici:
        foglie = 0
        for v in range(N):
            if (vertici >> v) & 1 and bin(ADJ[v] & vertici).count("1") <= 1:
                foglie |= 1 << v
        if not foglie:
            return False          # tutti i vertici rimasti hanno grado ≥ 2: contengono un ciclo
        vertici &= ~foglie
    return True


if __name__ == "__main__":
    inizio = time.time()
    n_insiemi = coppie = aciclici = 0
    for I in insiemi_indipendenti(12):
        n_insiemi += 1
        for u in range(N):
            if not (I >> u) & 1:
                coppie += 1
                if e_foresta(TUTTI & ~I & ~(1 << u)):
                    aciclici += 1
    print(f"insiemi indipendenti di taglia 12: {n_insiemi}; coppie (I,u): {coppie}; complementi aciclici: {aciclici}; "
          f"tempo {time.time() - inizio:.2f} s")
    print("(★) è", "VERA" if aciclici else "FALSA")

\end{lstlisting}
\begin{lstlisting}[caption={certificati/verifica\_dfs\_88.py}]
"""Verifica indipendente dell'etichettatura a 88 cammini: enumerazione esplicita per DFS.

Metodo diverso dalla ricorsione p(v): da ogni valle si esplorano tutti i cammini con etichette
crescenti e si contano uno a uno. Controlla anche che l'input sia una biiezione.
"""
ORDINE = ("00000,11110,11001,10101,01101,10011,01011,00111,10000,01000,00100,11100,00010,11010,"
          "10110,01110,00001,11111,11000,10100,01100,10010,01010,00110,10001,01001,00101,11101,"
          "00011,11011,10111,01111").split(",")
D = 5


def vicini(s):
    """Vicini della stringa s: cambia un carattere."""
    return [s[:i] + ("1" if s[i] == "0" else "0") + s[i + 1:] for i in range(D)]


def conta_da(s, etichetta):
    """Numero di cammini in salita che iniziano in s (s incluso come cammino di lunghezza 1)."""
    return 1 + sum(conta_da(w, etichetta) for w in vicini(s) if etichetta[w] > etichetta[s])


if __name__ == "__main__":
    assert len(set(ORDINE)) == 32 and all(len(s) == 5 and set(s) <= {"0", "1"} for s in ORDINE)
    etichetta = {s: i + 1 for i, s in enumerate(ORDINE)}
    valli = [s for s in ORDINE if all(etichetta[w] > etichetta[s] for w in vicini(s))]
    totale = sum(conta_da(s, etichetta) for s in valli)
    print("valli:", valli)
    print("cammini in salita (DFS esplicita):", totale)

\end{lstlisting}
\begin{lstlisting}[caption={certificati/verifica\_etichettatura\_q3.py}]
"""
Verifica indipendente della singola etichettatura consegnata per Q_3:
stampa valli, p(v) per ogni vertice e il totale dei cammini in salita (interi esatti).
"""
D = 3
CONSEGNA = ['000', '100', '010', '110', '101', '011', '001', '111']   # ordine crescente di etichetta


def vicini(s):
    """Stringhe 0/1 che differiscono da s in esattamente una coordinata."""
    return [s[:i] + ('1' if s[i] == '0' else '0') + s[i + 1:] for i in range(D)]


def verifica(consegna):
    """Calcola p(v) in ordine di etichetta e restituisce (valli, p, totale)."""
    assert sorted(consegna) == sorted(format(k, f'0{D}b') for k in range(2 ** D)), "non e' una biiezione"
    etichetta = {s: k + 1 for k, s in enumerate(consegna)}
    p, valli = {}, []
    for s in consegna:
        minori = [w for w in vicini(s) if etichetta[w] < etichetta[s]]
        if not minori:
            valli.append(s)
        p[s] = (1 if not minori else 0) + sum(p[w] for w in minori)
    return valli, p, sum(p.values())


if __name__ == "__main__":
    valli, p, totale = verifica(CONSEGNA)
    print("valli:", valli)
    for s in CONSEGNA:
        print(f"  {s}  etichetta {CONSEGNA.index(s)+1}  p = {p[s]}")
    print("cammini in salita totali:", totale)

\end{lstlisting}
\begin{lstlisting}[caption={certificati/verifica\_etichettatura\_q4.py}]
"""Verifica ESATTA (interi Python) dell'etichettatura candidata di Q_4 con 34 cammini in salita.
Due metodi indipendenti: (i) DP di Lemma 1 in ordine di etichetta; (ii) enumerazione esplicita
(DFS) di tutti i cammini in salita da ogni valle. Stampa la tabella p(v) per la consegna."""
D = 4
N = 1 << D
# Convenzione: carattere i-esimo della stringa (da sinistra) = coordinata i; bit i dell'intero.
ETICHETTATURA = ['0000', '1111', '0111', '1101', '1110', '1011', '0001', '1000',
                 '0010', '0100', '1001', '0110', '0011', '1010', '0101', '1100']


def da_stringa(s):
    """Converte la stringa 0/1 (carattere i = coordinata i) nell'intero con bit i = coordinata i."""
    return sum(int(c) << i for i, c in enumerate(s))


def vicini(v):
    """I d vicini di v in Q_d: si cambia esattamente una coordinata."""
    return [v ^ (1 << i) for i in range(D)]


def conta_con_dp(ordine, etichetta):
    """Metodo (i): p(v) = [v valle] + somma p(u) sui vicini u con etichetta minore."""
    p = {}
    righe = []
    for v in ordine:
        minori = [u for u in vicini(v) if etichetta[u] < etichetta[v]]
        p[v] = 1 if not minori else sum(p[u] for u in minori)
        righe.append((format(v, f'0{D}b')[::-1], etichetta[v], [format(u, f'0{D}b')[::-1] for u in minori], p[v]))
    return sum(p.values()), righe


def conta_con_dfs(etichetta):
    """Metodo (ii): enumera esplicitamente ogni cammino in salita partendo da ogni valle."""
    def estendi(v):
        totale = 1  # il cammino che termina qui
        for w in vicini(v):
            if etichetta[w] > etichetta[v]:
                totale += estendi(w)
        return totale
    valli = [v for v in range(N) if all(etichetta[w] > etichetta[v] for w in vicini(v))]
    return sum(estendi(v) for v in valli), valli


def main():
    ordine = [da_stringa(s) for s in ETICHETTATURA]
    assert sorted(ordine) == list(range(N)), "non e' una biiezione"
    etichetta = {v: i + 1 for i, v in enumerate(ordine)}
    totale_dp, righe = conta_con_dp(ordine, etichetta)
    totale_dfs, valli = conta_con_dfs(etichetta)
    for stringa, lab, minori, p in righe:
        print(f"{stringa}  etichetta {lab:2d}  vicini minori {minori}  p = {p}")
    print("valli:", [format(v, f'0{D}b')[::-1] for v in valli])
    print("totale (DP di Lemma 1):", totale_dp)
    print("totale (DFS esplicita):", totale_dfs)
    assert totale_dp == totale_dfs == 34


if __name__ == "__main__":
    main()

\end{lstlisting}
\begin{lstlisting}[caption={certificati/verifica\_q5\_bipartita.py}]
"""Certificato per il bound inferiore U(Q_5) >= 88 (implementazione 2, metodo diverso).

Stesso enunciato di verifica_q5_indipendenti.py, ma senza usare la simmetria:
ogni insieme indipendente e' A u B con A sottoinsieme dei vertici di peso pari e
B sottoinsieme dei vertici di peso dispari NON adiacenti ad A. Si scorrono tutti i 2^16
sottoinsiemi A (bitmask), poi tutti i B di taglia 12-|A| tra i dispari disponibili.
Aciclicita' con union-find (metodo diverso dal conteggio componenti). Rigore = exact.
"""
import itertools
import time

D = 5
N = 1 << D
PARI = [v for v in range(N) if bin(v).count("1") % 2 == 0]
DISPARI = [v for v in range(N) if bin(v).count("1") % 2 == 1]
SPIGOLI = [(v, v ^ (1 << i)) for v in range(N) for i in range(D) if v < v ^ (1 << i)]


def maschera_vicini(v):
    """Bitmask dei vicini di v."""
    return sum(1 << (v ^ (1 << i)) for i in range(D))


VIC = [maschera_vicini(v) for v in range(N)]


def ha_ciclo(maschera_vertici):
    """True sse il sottografo indotto dai vertici nella bitmask contiene un ciclo (union-find)."""
    padre = list(range(N))

    def trova(x):
        while padre[x] != x:
            padre[x] = padre[padre[x]]
            x = padre[x]
        return x

    for a, b in SPIGOLI:
        if (maschera_vertici >> a) & 1 and (maschera_vertici >> b) & 1:
            ra, rb = trova(a), trova(b)
            if ra == rb:
                return True
            padre[ra] = rb
    return False


def insiemi_indipendenti_12():
    """Genera le bitmask di tutti gli insiemi indipendenti di taglia 12."""
    for bits in range(1 << 16):
        A = [PARI[i] for i in range(16) if (bits >> i) & 1]
        if len(A) > 12:
            continue
        vietati = 0
        for a in A:
            vietati |= VIC[a]
        disponibili = [w for w in DISPARI if not (vietati >> w) & 1]
        for B in itertools.combinations(disponibili, 12 - len(A)):
            m = 0
            for x in A + list(B):
                m |= 1 << x
            yield m


if __name__ == "__main__":
    inizio = time.time()
    tutti = (1 << N) - 1
    n_insiemi, n_test, aciclici = 0, 0, 0
    for I in insiemi_indipendenti_12():
        n_insiemi += 1
        for u in range(N):
            if (I >> u) & 1:
                continue
            n_test += 1
            if not ha_ciclo(tutti & ~I & ~(1 << u)):
                aciclici += 1
    print(f"insiemi indipendenti di taglia 12 (tutti, senza simmetria): {n_insiemi}")
    print(f"coppie (I,u) esaminate: {n_test}; complementi aciclici: {aciclici}")
    print(f"tempo {time.time()-inizio:.2f}s")

\end{lstlisting}
\begin{lstlisting}[caption={certificati/verifica\_q5\_indipendenti.py}]
"""Certificato per il bound inferiore U(Q_5) >= 88 (implementazione 1: backtracking).

Enunciato verificato (Lemma computazionale):
  per ogni insieme indipendente I di Q_5 con |I| = 12 e per ogni vertice u di Q_5,
  il sottografo indotto da V meno (I u {u}) contiene un ciclo.
Riduzione: le traslazioni x -> x XOR t sono automorfismi di Q_5, quindi basta considerare
gli I che contengono il vertice 0 (ogni I non vuoto si trasla a contenerne uno).
Aciclicita' decisa con conteggio esatto: componenti + spigoli (un grafo e' una foresta sse
spigoli = vertici - componenti). Tutto intero: rigore = exact.
"""
import time
from conta_cammini import vicini

D = 5
N = 1 << D
ADJ = [set(vicini(v, D)) for v in range(N)]


def e_foresta(S):
    """True sse il sottografo indotto da S e' aciclico (spigoli = |S| - componenti)."""
    S = set(S)
    spigoli = sum(len(ADJ[v] & S) for v in S) // 2
    visti, comp = set(), 0
    for s in S:
        if s in visti:
            continue
        comp += 1
        pila, visti = [s], visti | {s}
        while pila:
            v = pila.pop()
            for w in ADJ[v] & S:
                if w not in visti:
                    visti.add(w)
                    pila.append(w)
    return spigoli == len(S) - comp


def indipendenti_con_zero(k):
    """Genera tutti gli insiemi indipendenti di taglia k contenenti 0 (vertici in ordine crescente)."""
    def ric(I, prossimo, vietati):
        if len(I) == k:
            yield list(I)
            return
        for v in range(prossimo, N):
            if v not in vietati and N - v >= k - len(I):
                yield from ric(I + [v], v + 1, vietati | ADJ[v] | {v})
    yield from ric([0], 1, ADJ[0] | {0})


if __name__ == "__main__":
    inizio = time.time()
    n_insiemi, n_test, violazioni = 0, 0, []
    for I in indipendenti_con_zero(12):
        n_insiemi += 1
        for u in range(N):
            if u in I:
                continue
            n_test += 1
            if e_foresta(set(range(N)) - set(I) - {u}):
                violazioni.append((I, u))
    print(f"insiemi indipendenti di taglia 12 contenenti 0: {n_insiemi}")
    print(f"coppie (I,u) esaminate: {n_test}; complementi aciclici trovati: {len(violazioni)}")
    print(f"tempo {time.time()-inizio:.2f}s")

\end{lstlisting}

\subsection{Letteratura arXiv consultata (ricerca deterministica, abstract letti)}
\begin{itemize}
\item \cite{1412.3893v1} The competition between simple and complex evolutionary trajectories in asexual populations — Ian E. Ochs, Michael M. Desai (2014). \emph{Query}: \texttt{uphill paths}.
\end{itemize}
\section{Problema 3 — Bulgarian solitaire}
\subsection{Enunciato ufficiale: preambolo e regole di consegna}
\subsubsection{Problema 3 --- Bulgarian
solitaire}\label{problema-3-bulgarian-solitaire}

\emph{(Enunciato formattato in LaTeX, contenuto fedele all'originale;
ricevuto il 2026-09-26.)}

\paragraph{Preambolo}\label{preambolo}

\textbf{Bulgarian solitaire.} Repeatedly take one card from every pile
and form a new pile. How many moves can it take before the piles start
to cycle?

A partition of a positive integer \(n\) is a weakly decreasing sequence
\(\lambda = (\lambda_1, \ldots, \lambda_s)\) of positive integers with
\(\lambda_1 + \cdots + \lambda_s = n\). Think of it as a division of
\(n\) cards into \(s\) piles.

The shift \(B(\lambda)\) is the partition of \(n\) obtained as follows:
remove one card from every pile, discard the piles that become empty,
and add one new pile consisting of the \(s\) removed cards. Formally,
\(B(\lambda)\) is the partition whose parts are the positive numbers
among \(\lambda_1 - 1, \ldots, \lambda_s - 1\) together with one extra
part equal to \(s\).

Iterating \(B\) from any starting partition must eventually repeat,
since there are finitely many partitions of \(n\). Call \(\lambda\)
\emph{cyclic} if \(B^i(\lambda) = \lambda\) for some \(i \ge 1\), and
let \[
d_B(\lambda) = \min\{\, i \ge 0 : B^i(\lambda) \text{ is cyclic} \,\}, \qquad
D_B(n) = \max\{\, d_B(\lambda) : \lambda \text{ is a partition of } n \,\}.
\] So \(D_B(n)\) is the largest number of shifts that can be needed
before the process starts to cycle.

Write \(T_k = k(k+1)/2\) and \(\delta_k = (k, k-1, \ldots, 2, 1)\), the
partition of \(T_k\) into distinct parts. Every \(n \ge 1\) satisfies
\(T_{k-1} < n \le T_k\) for exactly one \(k\), which we call the
\emph{rank} of \(n\).

Examples: \(B((2,1,1,1,1)) = (5,1)\); \(B((5,1)) = (4,2)\);
\(B((4,2)) = (3,2,1)\); \(B((3,2,1)) = (3,2,1)\). Here
\(d_B((2,1,1,1,1)) = 3\).

\textbf{What to hand in.} For each cell you attempt, hand in a written
proof. You may use a computer to explore. A proof may rely on a
computation only if you include the code, it runs in under 10 minutes on
a laptop, and the computation is exhaustive over a finite set that your
argument has reduced the problem to. Where a cell asks you to determine
a quantity, give the exact value as a formula in \(k\) and prove both
bounds; state any partitions you use explicitly as functions of \(k\).
Say clearly which cells you consider solved and which are partial.

Citing a published result for the statement you are asked to prove does
not count as a solution; a proof written out in full does, whatever its
source.

\subsection{Cella 1 (1 punti) — stato: in analisi}
\subparagraph{Parte 1 (C1) --- Cyclic partitions and
cycles}\label{parte-1-c1-cyclic-partitions-and-cycles}

\textbf{Punteggio:} 1 point · \textbf{Valutazione:} Judged

First, the long-run behaviour: which partitions repeat under the shift,
and how they fall into cycles. Let \(n = T_k\). Prove that for every
partition \(\lambda\) of \(n\) there is an \(i\) with
\(B^i(\lambda) = \delta_k\), and that \(\delta_k\) is the only cyclic
partition of \(n\). Then let \(n\) be arbitrary of rank \(k\), say
\(n = T_{k-1} + r\) with \(1 \le r \le k\): determine all cyclic
partitions of \(n\), and determine the number of distinct cycles of
\(B\) on the partitions of \(n\). Prove both.

\textbf{Enunciato protetto dato agli agenti (state.json).} \texttt{Let n = T\_k. Prove that every partition of n reaches delta\_k = (k,k-1,...,1) under iterated B and that delta\_k is the only cyclic partition. Then for n = T\_\{k-1\} + r, 1 \textless{}= r \textless{}= k, determine all cyclic partitions of n and the number of distinct cycles of B on partitions of n, with proofs.}

\textbf{Evidenza registrata in STATUS.md.} struttura nota (collane), prova da scrivere

\subsubsection{Consegna (prova)}
Nessuna bozza di consegna ancora scritta.

\subsubsection{Decisioni degli agenti}
\paragraph{Run \texttt{p3\_c1}}

\subsubsection{Codice su cui poggia il tentativo}
Nessun codice nei tentativi.

\subsection{Cella 2 (2 punti) — stato: in analisi}
\subparagraph{\texorpdfstring{Parte 2 (C2) --- \(D_B\) at triangular
\(n\)}{Parte 2 (C2) --- D\_B at triangular n}}\label{parte-2-c2-d_b-at-triangular-n}

\textbf{Punteggio:} 2 points · \textbf{Valutazione:} Judged

Next, how long the process can take to reach a cycle, starting with the
triangular numbers. Determine \(D_B(T_k)\) for every \(k\), with proof
of both bounds.

\textbf{Enunciato protetto dato agli agenti (state.json).} \texttt{Determine D\_B(T\_k) for every k, with proof of both bounds (exact value as a formula in k).}

\textbf{Evidenza registrata in STATUS.md.} atteso \$k\textasciicircum{}2-k\$ (Igusa/Etienne, da verificare)

\subsubsection{Consegna (prova)}
Nessuna bozza di consegna ancora scritta.

\subsubsection{Decisioni degli agenti}
\paragraph{Run \texttt{p3\_c2}}

\subsubsection{Codice su cui poggia il tentativo}
Nessun codice nei tentativi.

\subsection{Cella 3 (3 punti) — stato: in analisi}
\subparagraph{Parte 3 (C3) --- A general upper
bound}\label{parte-3-c3-a-general-upper-bound}

\textbf{Punteggio:} 3 points · \textbf{Valutazione:} Judged

Now the numbers strictly between two consecutive triangular numbers.
Prove that for every \(k \ge 4\) and every non-triangular \(n\) with
\(T_{k-1} < n < T_k\), \[
D_B(n) \le k^2 - 2k - 1,
\] and determine \(D_B(T_k - 1)\) exactly. Determine also, for that
\(n\), which partitions attain the maximum.

\textbf{Enunciato protetto dato agli agenti (state.json).} \texttt{Prove that for every k \textgreater{}= 4 and every non-triangular n with T\_\{k-1\} \textless{} n \textless{} T\_k, D\_B(n) \textless{}= k\textasciicircum{}2 - 2k - 1; determine D\_B(T\_k - 1) exactly and which partitions attain the maximum.}

\textbf{Evidenza registrata in STATUS.md.} —

\subsubsection{Consegna (prova)}
Nessuna bozza di consegna ancora scritta.

\subsubsection{Decisioni degli agenti}
\paragraph{Run \texttt{p3\_c3}}

\subsubsection{Codice su cui poggia il tentativo}
Nessun codice nei tentativi.

\subsection{Cella 4 (5 punti) — stato: in analisi}
\subparagraph{Parte 4 (C4) --- One above a triangular
number}\label{parte-4-c4-one-above-a-triangular-number}

\textbf{Punteggio:} 5 points · \textbf{Valutazione:} Judged

The first family just above a triangular number: \(n = T_{k-1} + 1\),
that is \(n = 11, 16, 22, 29, \ldots\) for \(k = 5, 6, 7, 8, \ldots\).
Determine \(D_B(T_{k-1} + 1)\) for every \(k \ge 5\), with proof of both
bounds.

\textbf{Enunciato protetto dato agli agenti (state.json).} \texttt{Determine D\_B(T\_\{k-1\} + 1) for every k \textgreater{}= 5, with proof of both bounds.}

\textbf{Evidenza registrata in STATUS.md.} —

\subsubsection{Consegna (prova)}
Nessuna bozza di consegna ancora scritta.

\subsubsection{Decisioni degli agenti}
\paragraph{Run \texttt{p3\_c4}}

\subsubsection{Codice su cui poggia il tentativo}
Nessun codice nei tentativi.

\subsection{Cella 5 (8 punti) — stato: in analisi}
\subparagraph{Parte 5 (C5) --- Two above a triangular
number}\label{parte-5-c5-two-above-a-triangular-number}

\textbf{Punteggio:} 8 points · \textbf{Valutazione:} Judged

The next family: \(n = T_{k-1} + 2\), that is
\(n = 3, 5, 8, 12, 17, 23, \ldots\) for
\(k = 2, 3, 4, 5, 6, 7, \ldots\). Determine \(D_B(T_{k-1} + 2)\) for
every \(k\), with proof of both bounds. State exactly for which \(k\)
your formula holds, and give the remaining values separately. Your upper
bound must be a single argument valid for all \(k\) in that range, not a
separate treatment of each \(k\), and you must give the extremal
partitions explicitly as a function of \(k\).

Say also where the straightforward extension of the C3 argument stops:
give the bound it does yield, show it is strictly weaker than the truth,
and identify precisely what your proof supplies in its place.

\textbf{Enunciato protetto dato agli agenti (state.json).} \texttt{Determine D\_B(T\_\{k-1\} + 2) for every k, with proof of both bounds; state exactly for which k the formula holds and give the remaining values; the upper bound must be one uniform argument; give the extremal partitions explicitly in k; say where the straightforward extension of the C3 argument stops and what replaces it.}

\textbf{Evidenza registrata in STATUS.md.} —

\subsubsection{Consegna (prova)}
Nessuna bozza di consegna ancora scritta.

\subsubsection{Decisioni degli agenti}
\paragraph{Run \texttt{p3\_c5}}

\subsubsection{Codice su cui poggia il tentativo}
Nessun codice nei tentativi.

\subsection{Cella 6 (13 punti) — stato: in analisi}
\subparagraph{\texorpdfstring{Parte 6 (C6) --- \(D_B(n)\) for every
\(n\)}{Parte 6 (C6) --- D\_B(n) for every n}}\label{parte-6-c6-d_bn-for-every-n}

\textbf{Punteggio:} 13 points · \textbf{Valutazione:} Judged ·
\textbf{Open question}

Finally, the whole function \(n \mapsto D_B(n)\). Determine \(D_B(n)\)
for every \(n\).

\emph{(Fine del problema 3. Punteggi: 1+2+3+5+8+13 = 32.)}

\textbf{Enunciato protetto dato agli agenti (state.json).} \texttt{OPEN. Determine D\_B(n) for every n.}

\textbf{Evidenza registrata in STATUS.md.} —

\subsubsection{Consegna (prova)}
Nessuna bozza di consegna ancora scritta.

\subsubsection{Decisioni degli agenti}
\paragraph{Run \texttt{p3\_c6}}

\subsubsection{Codice su cui poggia il tentativo}
Nessun codice nei tentativi.

\subsection{Letteratura arXiv consultata (ricerca deterministica, abstract letti)}
\begin{itemize}
\item \cite{math/0401385v2} Random Bulgarian solitaire — Serguei Popov (2004). \emph{Query}: \texttt{Bulgarian solitaire}.
\item \cite{1503.00885v1} The Bulgarian solitaire and the mathematics around it — Vesselin Drensky (2015). \emph{Query}: \texttt{Bulgarian solitaire}.
\item \cite{2607.17194v1} A short survey the game Bulgarian solitaire and related games — Romeo Meštrović (2026). \emph{Query}: \texttt{Bulgarian solitaire}.
\item \cite{1101.1546v3} Revisiting Toom's proof of Bulgarian Solitaire — Therese A. Hart, Gabriel Khan, Mizan R. Khan (2011). \emph{Query}: \texttt{Bulgarian solitaire}.
\item \cite{1703.07102v1} An exponential limit shape of random \$q\$-proportion Bulgarian solitaire — Kimmo Eriksson, Markus Jonsson abd Jonas Sjöstrand (2017). \emph{Query}: \texttt{Bulgarian solitaire}.
\item \cite{2208.14496v1} Limiting behavior in growth of Bulgarian Solitaire orbits — Nhung Pham (2022). \emph{Query}: \texttt{Bulgarian solitaire}.
\end{itemize}
\section{Problema 4 — Disjoint congruence classes and large gcd}
\subsection{Enunciato ufficiale: preambolo e regole di consegna}
\subsubsection{Problema 4 --- Disjoint congruence classes and large
gcd}\label{problema-4-disjoint-congruence-classes-and-large-gcd}

\emph{(Enunciato formattato in LaTeX, contenuto fedele all'originale;
ricevuto il 2026-09-26. Titolo della colonna non riportato nel testo
incollato; la domanda di apertura è: ``If congruence classes are
pairwise disjoint, must two of their moduli share a large common
factor?'')}

\paragraph{Preambolo}\label{preambolo}

For integers \(a\) and \(m \ge 1\) write \[
a \pmod m = \{\, a + mt : t \in \mathbb{Z} \,\}
\] for the congruence class of \(a\) modulo \(m\). A finite family of
classes \(a_1 \pmod{m_1}, \ldots, a_k \pmod{m_k}\) is \emph{pairwise
disjoint} if \(a_i \pmod{m_i} \cap a_j \pmod{m_j} = \emptyset\) whenever
\(i < j\). Nothing is assumed about the moduli beyond \(m_i \ge 1\):
they may repeat, and the residues \(a_i\) may repeat as well.

By the Chinese remainder theorem, two classes meet if and only if the
congruences \(x \equiv a_i \pmod{m_i}\) and \(x \equiv a_j \pmod{m_j}\)
are compatible, that is \[
a_i \pmod{m_i} \cap a_j \pmod{m_j} \ne \emptyset \iff \gcd(m_i, m_j) \mid a_i - a_j. \qquad (*)
\] The question of this column is how large the moduli of a disjoint
family must be forced to overlap in the gcd sense.

\textbf{Question.} Let \(k \ge 2\) and let
\(a_1 \pmod{m_1}, \ldots, a_k \pmod{m_k}\) be pairwise disjoint. Must
there be a pair \(i < j\) with \[
\gcd(m_i, m_j) \ge k \,?
\] This is believed to be true for every \(k\). It is sharp: the \(k\)
classes \(1, 2, \ldots, k \pmod k\) are pairwise disjoint --- their
differences have absolute value less than \(k\), so \(k\) divides none
of them --- and every pairwise gcd equals exactly \(k\). So \(k\) could
not be replaced by \(k+1\) in the statement.

\textbf{Examples.} The three classes \(0 \pmod 2\), \(1 \pmod 4\),
\(3 \pmod 8\) are pairwise disjoint by \((*)\), and their pairwise gcds
are \(2, 2, 4\); the largest is \(4 \ge 3\), as the question predicts
for \(k = 3\). By contrast \(0 \pmod 2\) and \(0 \pmod 3\) have coprime
moduli and do meet, at \(0\).

\textbf{Trivial bounds.} Two observations are immediate and you should
assume they are already on the table; presenting either as progress
counts for nothing.

\begin{enumerate}
\def\labelenumi{\arabic{enumi}.}
\tightlist
\item
  If \(\gcd(m_i, m_j) = 1\) then by \((*)\) the two classes meet. Hence
  in any pairwise disjoint family every pairwise gcd is at least \(2\).
  This answers the question for \(k = 2\), and for a family of size
  \(k\) that fails the question it confines every pairwise gcd to the
  range \(2 \le \gcd(m_i, m_j) \le k - 1\).
\item
  The class \(a \pmod m\) has density \(1/m\), so a pairwise disjoint
  family satisfies \(\sum_{i=1}^k 1/m_i \le 1\). In particular some
  modulus satisfies \(m_i \ge k\).
\end{enumerate}

Observation 2 bounds a single modulus from below; it says nothing about
any gcd, and the two together settle no size \(k \ge 3\). Every cell
below needs a genuinely new argument.

\textbf{What to hand in.} For each cell you attempt, hand in a written
proof. You may use a computer to explore. A proof may rely on a
computation only if you include the code, it runs in under ten minutes
on a laptop, and the computation is exhaustive over a finite set that
your own argument has reduced the problem to.

A computational cell is accepted only with an \emph{exhaustiveness
certificate}, meaning all four of:

\begin{enumerate}
\def\labelenumi{\arabic{enumi}.}
\tightlist
\item
  an exact description of the finite set your search enumerates,
  together with the proof that nothing outside it can be a
  counterexample;
\item
  every pruning rule you use, stated precisely, each with the proof that
  it discards only objects that cannot be completed to a counterexample;
\item
  the node count of the search at each size you claim, and the wall
  clock;
\item
  a rerun that reproduces those counts;
\item
  and, if anything at all survives your pruning at a size you claim, a
  separate decision for every survivor --- not a sample --- together
  with, for each one, the smallest part of it that already forces the
  decision, and your proof that no smaller part does. A survivor you
  cannot decide means the size is not certified.
\end{enumerate}

A run that reports ``I searched and found nothing'' without 1--5 is not
a certificate and scores nothing. If a search does not finish at some
size, report it as unfinished and do not claim that size; a partial
search is a legitimate partial result and should be handed in as one,
with the counts it reached.

Where a cell asks you to determine a boundary, state the exact value you
claim. Claim only what you have certified: an unfinished size, or one
certified under an assumption you have not verified, must be reported as
such. Say clearly which cells you consider solved and which are partial.

Citing a published result for the statement you are asked to prove does
not count as a solution; a proof written out in full does, whatever its
source. If you rely on a published argument, you are responsible for it:
check it yourself, and report it if it does not hold up.

\subsection{Cella 1 (1 punti) — stato: in analisi}
\subparagraph{Parte 1 (C1) --- Three
classes}\label{parte-1-c1-three-classes}

\textbf{Punteggio:} 1 point · \textbf{Valutazione:} Judged

The first size that the two trivial observations above do not settle.
Prove the statement for \(k = 3\): any three pairwise disjoint classes
have \(\gcd(m_i, m_j) \ge 3\) for some \(i < j\).

\textbf{Enunciato protetto dato agli agenti (state.json).} \texttt{Prove: any three pairwise disjoint congruence classes a\_i (mod m\_i) have gcd(m\_i,m\_j) \textgreater{}= 3 for some i \textless{} j.}

\textbf{Evidenza registrata in STATUS.md.} prova a mano abbozzata (parità)

\subsubsection{Consegna (prova)}
Nessuna bozza di consegna ancora scritta.

\subsubsection{Decisioni degli agenti}
\paragraph{Run \texttt{p4\_c1}}

\subsubsection{Codice su cui poggia il tentativo}
Nessun codice nei tentativi.

\subsection{Cella 2 (2 punti) — stato: in analisi}
\subparagraph{Parte 2 (C2) --- Four
classes}\label{parte-2-c2-four-classes}

\textbf{Punteggio:} 2 points · \textbf{Valutazione:} Judged

The same question for four classes. Prove the statement for \(k = 4\).

\textbf{Enunciato protetto dato agli agenti (state.json).} \texttt{Prove the statement for k = 4: any four pairwise disjoint classes have a pair with gcd(m\_i,m\_j) \textgreater{}= 4.}

\textbf{Evidenza registrata in STATUS.md.} schema "clique per primo" abbozzato

\subsubsection{Consegna (prova)}
Nessuna bozza di consegna ancora scritta.

\subsubsection{Decisioni degli agenti}
\paragraph{Run \texttt{p4\_c2}}

\subsubsection{Codice su cui poggia il tentativo}
Nessun codice nei tentativi.

\subsection{Cella 3 (3 punti) — stato: in analisi}
\subparagraph{\texorpdfstring{Parte 3 (C3) --- Every \(k\) up to
8}{Parte 3 (C3) --- Every k up to 8}}\label{parte-3-c3-every-k-up-to-8}

\textbf{Punteggio:} 3 points · \textbf{Valutazione:} Judged

A first range of sizes: after C1 and C2, the sizes \(k = 5, 6, 7\) and
\(8\) remain. Prove the statement for every \(k \le 8\).

\textbf{Enunciato protetto dato agli agenti (state.json).} \texttt{Prove the statement for every k \textless{}= 8 (sizes 5,6,7,8 remain after cells 1-2). A computation is admissible only with the full exhaustiveness certificate (finite set + proof, pruning rules proved, node counts, wall clock, rerun, decisions for every survivor).}

\textbf{Evidenza registrata in STATUS.md.} riduzione a insieme finito abbozzata (note.md)

\subsubsection{Consegna (prova)}
Nessuna bozza di consegna ancora scritta.

\subsubsection{Decisioni degli agenti}
\paragraph{Run \texttt{p4\_c3}}

\subsubsection{Codice su cui poggia il tentativo}
Nessun codice nei tentativi.

\subsection{Cella 4 (5 punti) — stato: in analisi}
\subparagraph{\texorpdfstring{Parte 4 (C4) --- Every \(k\) up to
12}{Parte 4 (C4) --- Every k up to 12}}\label{parte-4-c4-every-k-up-to-12}

\textbf{Punteggio:} 5 points · \textbf{Valutazione:} Judged

A longer range, and a precise account of your method at the sizes just
beyond it. Prove the statement for every \(k \le 12\). Then, for each
\(k\) from \(9\) to \(16\) in turn, report exactly what your method
leaves undecided at that \(k\): if nothing, say so and prove it; if
something, exhibit it in full. If that disagrees with any source you
consulted, say which of the two is right and why. An answer that reports
agreement with a source it did not test will be marked wrong.

\textbf{Enunciato protetto dato agli agenti (state.json).} \texttt{Prove the statement for every k \textless{}= 12; then for each k = 9..16 report exactly what the method leaves undecided (nothing: prove it; something: exhibit it in full); compare with sources only if tested.}

\textbf{Evidenza registrata in STATUS.md.} —

\subsubsection{Consegna (prova)}
Nessuna bozza di consegna ancora scritta.

\subsubsection{Decisioni degli agenti}
\paragraph{Run \texttt{p4\_c4}}

\subsubsection{Codice su cui poggia il tentativo}
Nessun codice nei tentativi.

\subsection{Cella 5 (8 punti) — stato: in analisi}
\subparagraph{Parte 5 (C5) --- The certified
boundary}\label{parte-5-c5-the-certified-boundary}

\textbf{Punteggio:} 8 points · \textbf{Valutazione:} Judged

The main certification cell: an initial range of sizes as long as you
can make it, plus two isolated sizes further out. Determine the largest
\(k\) for which you can certify the statement for all sizes up to and
including \(k\), and certify it. The range you claim must be contiguous,
and if your argument at a given size assumes that all smaller sizes have
already been settled you must say so. Then decide the two isolated sizes
\(k = 24\) and \(k = 30\): show that neither is the least size at which
the statement can fail.

For this cell the exhaustiveness certificate of the hand-in rules is not
enough on its own. Add:

\begin{enumerate}
\def\labelenumi{(\roman{enumi})}
\item
  a demonstration that your method is not vacuous. Construct a case in
  which an admissible configuration genuinely exists, run your machinery
  on it unmodified, and show it returns that configuration. A method
  that reports ``nothing survives'' at every size it is pointed at is
  indistinguishable from a method with a bug, and will be graded as one;
\item
  for every object your pruning leaves undecided, at every \(k\) you
  claim --- not a sample --- an individual decision, plus the smallest
  part of that object which already forces the decision, plus a proof
  that no smaller part does;
\item
  the exact list, not merely the count, of what survives at each \(k\);
\item
  every pruning rule you use beyond those you have proved, proved. If
  your search needs a rule you invented to finish a size, that rule is
  part of your claim: state it, prove it, and show the survivor list is
  unchanged when you switch it off. A size that only closes with an
  unproved rule is not certified;
\item
  a second implementation, written independently of your first, that
  differs in method --- not the same algorithm twice --- and the two
  survivor lists compared elementwise at every \(k\) you claim. Report
  any difference rather than reconciling it silently: a disagreement
  means one of them is wrong, and finding which is part of the cell.
\end{enumerate}

\textbf{Enunciato protetto dato agli agenti (state.json).} \texttt{Determine the largest k for which the statement is certified for all sizes up to k (contiguous range), with the strengthened certificate (i)-(v): non-vacuity demonstration, individual decisions for survivors with minimal forcing part, exact survivor lists, all pruning rules proved and tested off, a second independent implementation compared elementwise. Also decide k = 24 and k = 30: show neither is the least size at which the statement can fail.}

\textbf{Evidenza registrata in STATUS.md.} richiede 2 implementazioni indipendenti

\subsubsection{Consegna (prova)}
Nessuna bozza di consegna ancora scritta.

\subsubsection{Decisioni degli agenti}
\paragraph{Run \texttt{p4\_c5}}

\subsubsection{Codice su cui poggia il tentativo}
Nessun codice nei tentativi.

\subsection{Cella 6 (13 punti) — stato: in analisi}
\subparagraph{Parte 6 (C6) --- Beyond the
boundary}\label{parte-6-c6-beyond-the-boundary}

\textbf{Punteggio:} 13 points · \textbf{Valutazione:} Judged ·
\textbf{Open question}

Three directions beyond the certified range; any one of them counts. Any
of the following.

\begin{enumerate}
\def\labelenumi{(\alph{enumi})}
\item
  Decide a size \(k \ge 25\).
\item
  The asymptotic form. It is known that a pairwise disjoint family of
  size \(k\) always has a pair with \[
  \gcd(m_i, m_j) \;\ge\; k \cdot \exp\!\left( -(2 + o(1)) \frac{\log k}{\log\log k} \right),
  \] which is \(k^{1 - o(1)}\) but not linear in \(k\). Prove the
  statement in full, or prove the weaker bound \(\gcd(m_i, m_j) \ge ck\)
  for some absolute constant \(c > 0\), or improve the exponential
  factor above.
\item
  The group form. Let \(G\) be a group, let \(G_1, \ldots, G_k\) be
  subgroups of finite index \(n_i = [G : G_i]\), and let
  \(x_1 G_1, \ldots, x_k G_k\) be pairwise disjoint cosets. Is there a
  pair \(i < j\) with \(\gcd(n_i, n_j) \ge k\)? This is known for
  \(k \le 5\) and open for every \(k \ge 6\); settling \(k = 6\) counts
  as progress.
\end{enumerate}

\emph{(Fine del problema 4. Punteggi: 1+2+3+5+8+13 = 32.)}

\textbf{Enunciato protetto dato agli agenti (state.json).} \texttt{OPEN. Any of: (a) decide a size k \textgreater{}= 25; (b) prove gcd \textgreater{}= ck for an absolute c \textgreater{} 0 or improve the known factor k exp(-(2+o(1)) log k / log log k); (c) the group form (disjoint cosets of finite-index subgroups), settle k = 6.}

\textbf{Evidenza registrata in STATUS.md.} —

\subsubsection{Consegna (prova)}
Nessuna bozza di consegna ancora scritta.

\subsubsection{Decisioni degli agenti}
\paragraph{Run \texttt{p4\_c6}}

\subsubsection{Codice su cui poggia il tentativo}
Nessun codice nei tentativi.

\subsection{Letteratura arXiv consultata (ricerca deterministica, abstract letti)}
\begin{itemize}
\item \cite{2603.26043v1} Finiteness of Disjoint Covering Systems with Precisely One Repeated Modulus — Yu Hashimoto (2026). \emph{Query}: \texttt{disjoint covering systems}.
\item \cite{1511.04293v1} Searching for Disjoint Covering Systems with Precisely One Repeated Modulus — Shalosh B. Ekhad, Aviezri S. Fraenkel, Doron Zeilberger (2015). \emph{Query}: \texttt{disjoint covering systems}.
\item \cite{2607.24655v1} On the problem of large gcd for disjoint residue classes — Jan Fornal, Yu-Chen Sun (2026). \emph{Query}: \texttt{disjoint residue classes}.
\item \cite{2608.15873v1} Two Questions on \$G\$-harmonic Tuples — Murali Menon (2026). \emph{Query}: \texttt{disjoint residue classes}.
\end{itemize}
\begin{thebibliography}{99}
\bibitem{1101.1546v3} Therese A. Hart, Gabriel Khan, Mizan R. Khan. \emph{Revisiting Toom's proof of Bulgarian Solitaire}. arXiv:1101.1546v3 (2011). \url{http://arxiv.org/abs/1101.1546v3}
\bibitem{1204.3850v1} Jérémie Chalopin, Shantanu Das, Yann Disser, Matúš Mihalák, Peter Widmayer. \emph{Simple Agents Learn to Find Their Way: An Introduction on Mapping Polygons}. arXiv:1204.3850v1 (2012). \url{http://arxiv.org/abs/1204.3850v1}
\bibitem{1412.3893v1} Ian E. Ochs, Michael M. Desai. \emph{The competition between simple and complex evolutionary trajectories in asexual populations}. arXiv:1412.3893v1 (2014). \url{http://arxiv.org/abs/1412.3893v1}
\bibitem{1503.00885v1} Vesselin Drensky. \emph{The Bulgarian solitaire and the mathematics around it}. arXiv:1503.00885v1 (2015). \url{http://arxiv.org/abs/1503.00885v1}
\bibitem{1511.04293v1} Shalosh B. Ekhad, Aviezri S. Fraenkel, Doron Zeilberger. \emph{Searching for Disjoint Covering Systems with Precisely One Repeated Modulus}. arXiv:1511.04293v1 (2015). \url{http://arxiv.org/abs/1511.04293v1}
\bibitem{1703.07102v1} Kimmo Eriksson, Markus Jonsson abd Jonas Sjöstrand. \emph{An exponential limit shape of random \$q\$-proportion Bulgarian solitaire}. arXiv:1703.07102v1 (2017). \url{http://arxiv.org/abs/1703.07102v1}
\bibitem{1801.07837v1} Dmitriy Bilyk, Ryan W Matzke. \emph{On the Fejes Tóth Problem about the Sum of Angles Between Lines}. arXiv:1801.07837v1 (2018). \url{http://arxiv.org/abs/1801.07837v1}
\bibitem{2007.08698v2} Tongseok Lim, Robert J. McCann. \emph{On Fejes Tóth's conjectured maximizer for the sum of angles between lines}. arXiv:2007.08698v2 (2020). \url{http://arxiv.org/abs/2007.08698v2}
\bibitem{2208.14496v1} Nhung Pham. \emph{Limiting behavior in growth of Bulgarian Solitaire orbits}. arXiv:2208.14496v1 (2022). \url{http://arxiv.org/abs/2208.14496v1}
\bibitem{2603.26043v1} Yu Hashimoto. \emph{Finiteness of Disjoint Covering Systems with Precisely One Repeated Modulus}. arXiv:2603.26043v1 (2026). \url{http://arxiv.org/abs/2603.26043v1}
\bibitem{2607.17194v1} Romeo Meštrović. \emph{A short survey the game Bulgarian solitaire and related games}. arXiv:2607.17194v1 (2026). \url{http://arxiv.org/abs/2607.17194v1}
\bibitem{2607.24655v1} Jan Fornal, Yu-Chen Sun. \emph{On the problem of large gcd for disjoint residue classes}. arXiv:2607.24655v1 (2026). \url{http://arxiv.org/abs/2607.24655v1}
\bibitem{2608.15873v1} Murali Menon. \emph{Two Questions on \$G\$-harmonic Tuples}. arXiv:2608.15873v1 (2026). \url{http://arxiv.org/abs/2608.15873v1}
\bibitem{math/0401385v2} Serguei Popov. \emph{Random Bulgarian solitaire}. arXiv:math/0401385v2 (2004). \url{http://arxiv.org/abs/math/0401385v2}
\bibitem{bm18} D. Bilyk, R. Matzke. \emph{On the Fejes Tóth problem about the sum of angles between lines}. Proc. AMS (2019); letta integralmente, vedi problema-1/fonti/.
\end{thebibliography}
\end{document}
